% IEEE TKDE layout simulation (article class, 10pt, twocolumn)
% For page count estimation. True compile: use Overleaf with IEEEtran.cls
% Content faithful to v77 (Paper3A_Draft_v77_TKDE.docx)
\documentclass[10pt,twocolumn,a4paper]{article}
\usepackage[a4paper,top=19mm,bottom=24mm,left=14.5mm,right=14.5mm,columnsep=5mm]{geometry}
\usepackage[T1]{fontenc}
\usepackage[utf8]{inputenc}
\usepackage{amsmath,amssymb,amsthm}
\usepackage{graphicx}
\usepackage{booktabs}
\usepackage{placeins}            % manual \FloatBarrier where needed
\usepackage{xcolor}
\usepackage{url}
\usepackage{microtype}
\usepackage{times}
\usepackage{cite}
\usepackage{balance}
\usepackage[hidelinks]{hyperref}
\usepackage{titlesec}
\titleformat{\section}{\bfseries\normalsize}{\Roman{section}.}{0.5em}{}
\titleformat{\subsection}{\bfseries\itshape\small}{\Alph{subsection}.}{0.4em}{}
\setlength\parindent{7pt}\setlength\parskip{0pt}
\titlespacing*{\section}{0pt}{4pt}{2pt}
\titlespacing*{\subsection}{0pt}{3pt}{1pt}
\newtheorem{theorem}{Theorem}
\newtheorem{lemma}{Lemma}
\newtheorem{definition}{Definition}
\newtheorem{corollary}{Corollary}
\newtheorem{remark}{Remark}
\newtheorem{proposition}{Proposition}
%% HSM shorthand commands
\newcommand{\SHSM}{S_{\mathrm{HSM}}}
\newcommand{\SR}{S_{R}}
\newcommand{\SV}{S_{V}}
\newcommand{\ST}{S_{T}}
\newcommand{\SA}{S_{A}}
\newcommand{\SPER}{S_{P}}
\newcommand{\Nstar}{N^{*}_{\mathrm{pts}}}
\newcommand{\Npts}{N_{\mathrm{pts}}}
\newcommand{\thetastar}{\theta^{*}}
\newcommand{\dHSM}{d_{\mathrm{HSM}}}
\newcommand{\TPR}{\mathrm{TPR}}
\newcommand{\TNR}{\mathrm{TNR}}
%% ---------------------------------------------------------------------------
%% HSM v2 backfill (2026-04-15)
%%   All TBD/TBDcell/REMARK macros previously used as placeholders have been
%%   resolved with measured values from the v2 sweep.  Source CSVs:
%%     results/{summary_all,throughput,overhead,scale_analysis,
%%              noise_sensitivity,window_sweep}.csv
%%     code/results/{sdss,job,oltp,burst,burst_v3}_validation/.
%%   The placeholder macros (\TBD, \TBDcell, \REMARK) are kept defined below
%%   so the document still compiles if a future revision needs to re-introduce
%%   them, but no live invocations remain in the body.
%% ---------------------------------------------------------------------------
\usepackage{soul}
\newcommand{\TBD}[2]{#2} % placeholder retired: render second arg only.
\newcommand{\TBDcell}[1]{--}
\newcommand{\REMARK}[1]{}
\begin{document}
\twocolumn[{\begin{center}
{\large\bfseries HSM: Workload-Similarity Gating for
Index-Maintenance Decisions, with Formal Bounds}\\[4pt]
{\normalsize Arun Reungsilpkolkarn}\\
{\small School of Information Technology and Innovation, Bangkok University, Pathum Thani, Thailand.
\texttt{arun.r@bu.ac.th}}\\[6pt]
\end{center}
\begin{quote}\small
\textbf{Abstract---}Existing index advisors (Dexter, Supabase
\texttt{index\_advisor}, BALANCE, Indexer++) answer \emph{which}
indexes to build but invoke their $\Omega(N\log N)$ search at
every window or on a fixed schedule.  We address the complementary
question---\emph{when} should the advisor be invoked?---with HSM
(Hierarchical Similarity Measurement), a training-free,
DBMS-agnostic decision layer that gates advisor invocation in
$\Theta(\Npts)$ time.  HSM quantifies workload similarity along
five dimensions (template rank, volume, type, attribute overlap,
DWT-based temporal pattern) and exposes a single calibrated
threshold $\thetastar(N,Q)$.  Six theorems and four lemmas
establish metric soundness, complexity, economic gating
optimality, detector-quality bounds, a deployment speedup
certificate, and a closed-form optimal window size.
On PostgreSQL\,16 with TPC-H (SF\,0.2--3.0), HSM attains
$J\!\ge\!0.96$ at every scale and reduces advisor invocations by
$80.8\%$ versus always-on and $42.5\%$ versus periodic
($K\!=\!3$), at a kernel cost of $\approx\!1$\,ms per decision.
End-to-end deployment with two production PostgreSQL advisors
(Dexter, Supabase \texttt{index\_advisor}) demonstrates $66$--$71\%$
wall-clock savings versus always-on with $100\%$ trigger precision.
Four external workloads (SDSS, JOB/IMDB, pgbench, burst) and
cross-engine validation on MongoDB~7 confirm the kernel transfers
unmodified; only $\theta$ requires per-engine recalibration---no
kernel retraining.\\[2pt]
\textbf{Index Terms---}Decision layer, index-maintenance timing,
proactive index management, workload similarity,
hierarchical similarity measurement, discrete wavelet transform,
self-tuning databases.
\end{quote}\vspace{4pt}}]

\section{Introduction}

Database index performance is sensitive to workload characteristics. When
the distribution of query types, access frequency, or temporal patterns
changes substantially, an index configuration that was optimal for a
previous workload window may become inefficient for the current one.
Proactive index management---deciding whether to retain an existing index
or rebuild it based on workload evolution---requires a reliable mechanism
to determine how similar the current workload window is to the reference
window under which the index was last calibrated. This paper addresses
that prerequisite.

Prior work~\cite{Reungsinkonkarn2026iEEECON,Reungsilpkolkarn2026ICICT}
validated the HSM concept, demonstrating 77\% classification accuracy and
a 4{,}623$\times$ latency improvement over unoptimised baselines using a
%% These two prior-work numerics (77\% and 4623x) are *not* placeholders
%% --- they refer to the published v1 conference papers and stay frozen.

coarse interval-overlap metric $S_i=|T_{\mathrm{overlap}}|/|T_{\mathrm{total}}|$.
That metric cannot distinguish burst events, periodic sub-patterns, or
multi-scale frequency structure, and those studies provided no formal
theoretical analysis.

We introduce the HSM framework, making five contributions.
\textbf{C1}~A five-dimensional similarity measure with formal metric
properties, capturing burst events, periodic oscillations (detected
via angular cosine for phase-alignment), and long-term trends
simultaneously.
\textbf{C2}~A theoretical core of six theorems and four lemmas
covering wavelet pattern preservation, minimal-sufficient dimensions,
economic gating optimality, stratified-Hoeffding detector-quality
bounds, a deployment speedup certificate, and a closed-form optimal
window size (Section~IV).
\textbf{C3}~Access-attribute independence: HSM incurs no
reconstruction cost when the accessed column set changes
(Theorem~\ref{thm:econ}).
\textbf{C4}~A closed-form threshold $\thetastar(N,Q)$ validated
across SF\,0.2--3 (Youden $J\!\ge\!0.96$ at every scale) and
generalising across TPC-H, burst, and cross-engine workloads.
\textbf{C5}~Cross-engine validation on MongoDB~7, confirming that
the HSM kernel is DBMS-agnostic and requires only per-engine
$\theta$ recalibration---no kernel retraining.

HSM is a measurement layer rather than a physical index: $\Theta(\Npts)$
per decision, dominated by the $\Theta(N\log N)$ reconstruction it
gates.

\section{Related Work}

\subsection{Workload Characterisation and Similarity Methods}

Early workload compression~\cite{Schnaitter2006COLT} and automated
index selection~\cite{Chaudhuri1997Index} treat workloads as static
query sets via predicate overlap or template matching.  Recent
advisors take richer views but address \emph{what} to index rather
than \emph{when}: BALANCE~\cite{Wang2024BALANCE} uses self-supervised
contrastive learning over index-selection instances;
MFIX~\cite{Chang2024MFIX} applies multi-fidelity Bayesian
optimisation; Indexer++~\cite{Sharma2022IndexerPP} detects binary
trend shifts with BERT embeddings and deep RL.  All assume the
advisor is always invoked, and most rely on GPUs and labelled
corpora that may be unavailable cross-engine.
ShapeShifter~\cite{Wang2025ShapeShifter} and the benchmark
of~\cite{Brucato2025WorkloadSim} use similarity only as a binary
trigger.  Prior HSM work~\cite{Reungsinkonkarn2026iEEECON,%
Reungsilpkolkarn2026ICICT} validated the concept empirically; the
present paper adds formal proofs and scale-sensitivity analysis.

\smallskip\noindent\emph{Gap~2.1: We are not aware of a prior method
that jointly provides a continuous multi-dimensional similarity score
with formal metric properties, a closed-form decision threshold, and a
training-free deployment model for online indexing.}

\smallskip\noindent\textbf{Deployed Index Advisors.}
Alongside academic prototypes, two open-source PostgreSQL advisors
define the practical deployment state of automatic index
recommendation: \emph{Dexter}~\cite{AnkaneDexter}, which scans query
logs and applies what-if analysis via
HypoPG~\cite{Giraud2015HypoPG}\footnote{HypoPG creates hypothetical
indexes whose presence can be tested through the PostgreSQL query
planner without the $\Theta(N\log N)$ cost of physical
construction.} to recommend $B$-tree indexes; and Supabase
\texttt{index\_advisor}~\cite{SupabaseIndexAdvisor}, a PostgreSQL
extension that drives the same HypoPG mechanism from an SQL
interface. Both are widely deployed
(Heroku/Instacart/Supabase production use), actively maintained
through 2024, and have been adopted as baseline advisors in several
recent industrial and academic studies. Open-source commercial
derivatives include \emph{pganalyze} Index Advisor and
Huawei's \emph{DBMind}~\cite{Zhou2021DBMind}; the latter applies a
reinforcement-learning policy in openGauss (a PostgreSQL fork) and
therefore lies outside the vanilla-PostgreSQL protocol used in our
experiments. Neither Dexter nor Supabase \texttt{index\_advisor}
addresses \emph{when} to invoke advisor analysis---both assume the
caller decides invocation externally---so HSM composes with them
orthogonally by supplying the missing timing layer
(Table~\ref{tab:advisor-applicability}).
Both advisors are deployed end-to-end with HSM as the gating layer
in \S\ref{sec:a13d} (A13d), with measured wall-clock, advisor cost,
and trigger precision reported in Table~\ref{tab:a13d-results}.

\subsection{Signal Processing for Time-Series Similarity}

DWT~\cite{Daubechies1992Wavelets} decomposes signals into
multi-resolution components; Chan \& Fu~\cite{Chan2003Wavelets} applied
this to time-series indexing. SAX~\cite{Lin2007SAX} converts normalised
series to discrete symbols with bounded error; a 2025 survey~\cite{Pappa2025SAXSurvey}
reviews variants, including Quantile SAX~\cite{Silveira2023QuantileSAX}
for non-Gaussian series. DTW~\cite{Sakoe1978DTW} measures elastic
similarity; FastDTW~\cite{Salvador2007FastDTW} approximates DTW in linear
time, enabling $O(N)$ workload-window comparison. The three are complementary---multi-resolution decomposition (DWT),
bounded-error symbolisation (SAX), and elastic alignment
(FastDTW)---motivating the three-stage $\SPER$ pipeline
(Section~III).

\textit{Theoretical foundations.} Metric space theory~\cite{Deza2009Distances}
formalises similarity properties. Information-theoretic bounds for
symbolic representations draw on Cover \& Thomas~\cite{Cover2006InfoTheory}.
Wavelet approximation bounds are established in Mallat~\cite{Mallat1989Wavelets}.
These classical results provide the mathematical scaffolding for Section~IV.

\smallskip\noindent\emph{Gap~2.2: No prior work provides formal guarantees
for multi-scale workload similarity measurement---neither proving metric
properties nor establishing parameter optimality bounds.}

\subsection{Database Optimisation and Access-Attribute Dependence}

Self-tuning systems (AI-DBA~\cite{Li2022AIDBA}, CDBTune~\cite{Zhang2019CDBTune},
IA$^2$~\cite{Xu2024IA2}, GPTuner~\cite{Lao2024GPTuner}) all depend on
workload characterisation as an upstream step; none provides a principled
multi-dimensional similarity score with formal properties.

All physical index methods share \emph{access-attribute dependence}:
the structure is built for a fixed column set, and changing that set
requires full reconstruction ($\Theta(N\log N)$). This holds from
classical B-Trees~\cite{Ramakrishnan2003DMS} and Database
Cracking~\cite{Idreos2011Cracking}, through first-generation
learned indexes~\cite{Kraska2018LearnedIndex} (ALEX~\cite{Ding2020ALEX},
PGM-Index~\cite{Ferragina2020PGM}, XIndex~\cite{Tang2020XIndex}),
to multi-dimensional extensions
(Flood~\cite{Nathan2020Flood}, Tsunami~\cite{Ding2020Tsunami},
WaZI~\cite{Pai2024WaZI}).
A formal comparison appears in the Supplementary
(Structural Comparison with Prior Methods).

\smallskip\noindent\emph{Gap~2.3: No existing framework---B-Trees,
cracking, learned indexes (PGM, ALEX, XIndex), or workload-aware
partitioning (WaZI)---formally characterises access-attribute
independence or provides multi-dimensional temporal similarity with
cost guarantees.}

\subsection{NoSQL and Document Store Indexing}

Cattell's taxonomy~\cite{Cattell2011NoSQL} established the SQL/NoSQL
divide as a data-model and consistency tradeoff, but did not address
workload similarity measurement. On the schema side,
NoSE~\cite{Mior2017NoSE} formulates NoSQL physical design as an integer
program over a conceptual data model; however, it targets Cassandra's
wide-column model and does not consider temporal workload drift.
Klettke et al.~\cite{Klettke2016SchemaEvol} address schema evolution
and data migration in document stores at scale, highlighting that
agile NoSQL development produces frequent schema versions---a setting
where proactive drift detection is especially valuable.
De~Espona~Pernas et al.~\cite{deEspona2023MongoIndex} propose
query-only automatic index suggestion for MongoDB aggregation
pipelines, the closest NoSQL-specific index advisor to date; their
method does not model temporal windows or workload similarity.
On the SQL side, industrial index managers such as
AIM~\cite{Yadav2023AIM} and workload-reduction techniques like
Wred~\cite{Brucato2024Wred} remain relational-only.
Stonebraker and Pavlo's recent retrospective~\cite{Stonebraker2024WGA}
observes that relational and document models increasingly coexist in
modern deployments, yet no cross-engine workload similarity framework
has been proposed.

\smallskip\noindent\emph{Gap~2.4: No existing index management
framework---relational or NoSQL---provides a DBMS-agnostic workload
similarity measure validated across both SQL and document-store
engines with formal metric guarantees.}

These four gaps map to contributions C1--C5 (Section~I):
Gaps~2.1--2.2 $\to$ C1--C2; Gap~2.3 $\to$ C3--C5; Gap~2.4 $\to$
cross-engine MongoDB validation (Section~\ref{sec:mongo-exp}).

\section{The HSM Framework}

HSM compares two workload windows $W_A, W_B$, each of duration $T$
containing $\Npts$ time-stamped query events, in four phases:
preprocessing, five-dimensional similarity computation, score
aggregation, and index decision.

\subsection{Workload Preprocessing}

Each window $W$ is reduced to: a $k$-element query-template frequency
vector $\mathbf{f}\!\in\!\mathbb{R}^k$ (one entry per distinct template,
extracted by the DBMS-specific extractor; Section~III-B); a QPS time
series $q(t)$ at 1-second resolution giving $\Npts$ points\footnote{The
implementation normalises $q(t)$ by $\max_t q(t)$ before the wavelet
decomposition (\texttt{hsm\_v2\_kernel.py}: \texttt{sp\_v2}). This makes
the bound constant $M\!=\!1$ in Theorem~\ref{thm:wavelet}, so the cited
tight bound $\varepsilon\!\le\!0.071$ at $L\!=\!3$, $\Npts\!=\!100$
holds without further scaling.}; relation set
$\mathcal{T}$ and column set $\mathcal{C}$ forming the access-attribute
schema $A(W)\!=\!(\mathcal{T},\mathcal{C})$; and total volume
$\mathrm{QPS}\!=\!\sum_i f_i/T$.

\subsection{Five Similarity Dimensions}

\textbf{Ratio Similarity $\SR$} measures the relative-frequency
distribution of query templates via Spearman rank correlation on the
normalised query-frequency rank vectors $\mathbf{r}$ (one entry per
distinct template, derived from $\mathbf{f}$):
\begin{equation}
  \SR(W_A,W_B) = 1 - \frac{\arccos\!\bigl(\rho_s(\mathbf{r}_A,\mathbf{r}_B)\bigr)}{\pi},
\end{equation}
where $\rho_s(\cdot)$ is the Spearman rank correlation.
Since $\rho_s\!\in\![-1,1]$, we have $\arccos(\rho_s)\!\in\![0,\pi]$,
so $\SR\!\in\![0,1]$. Importantly, $d_R\!=\!1\!-\!\SR\!=\!\arccos(\rho_s)/\pi$
is the normalised \emph{correlation angle} and satisfies the triangle
inequality analytically \cite{Deza2009Distances}.
$\SR\!=\!1$ iff both windows share identical query-frequency rank order.
\textbf{Volume Similarity $\SV$} uses log-scale to handle
order-of-magnitude differences:
\begin{equation}
  \SV(W_A,W_B) = \exp(-|\ln Q_A - \ln Q_B|) = \frac{\min(Q_A,Q_B)}{\max(Q_A,Q_B)}.
\end{equation}
\textbf{Type Similarity $\ST$} compares query-template distributions via
angular distance on unit-normalised frequency vectors.
Let $\mathcal{Q}\!=\!\{q_1,\ldots,q_k\}$ be the set of $k$ distinct
operation categories observed in the monitoring period
(e.g.\ SQL query templates $\{$Q1, Q3, \ldots$\}$ for relational DBMS;
pipeline categories $\{$\texttt{find}, \texttt{aggregate},
\texttt{insert}, \texttt{update}, \texttt{delete}$\}$ for document DBMS).
Define $\mathbf{v}\!\in\!\mathbb{R}^k$ as the normalised frequency count
over $\mathcal{Q}$ within window $W$:
\begin{equation}
  \ST(W_A,W_B) = 1 - \frac{2}{\pi}\arccos\!\bigl(\hat{\mathbf{v}}_A\cdot\hat{\mathbf{v}}_B\bigr),
\end{equation}
where $\hat{\mathbf{v}}$ denotes the $\ell_2$-normalised unit vector.
Frequency counts are non-negative, so
$\hat{\mathbf{v}}_A\cdot\hat{\mathbf{v}}_B\!\in\![0,1]$,
giving $\ST\!\in\![0,1]$.
The dissimilarity $d_T\!=\!(2/\pi)\arccos(\hat{\mathbf{v}}_A\cdot\hat{\mathbf{v}}_B)$
is the \emph{angular distance} (geodesic) on $\mathbb{S}^{k-1}$,
satisfying the triangle inequality for any $k\!\geq\!1$
\cite{Deza2009Distances}.
This DBMS-agnostic formulation allows $\mathcal{Q}$ to be defined
by a \emph{DBMS-specific extractor} (Section~III-B), making $\ST$
applicable across heterogeneous database paradigms.

\textbf{Access Similarity $\SA$} measures schema-level overlap as an
equal-weight average of two Jaccard scores over table-level and
column-level access sets:
\begin{equation}
  \SA = \tfrac{1}{2}J(\mathcal{T}_A,\mathcal{T}_B)
       + \tfrac{1}{2}J(\mathcal{C}_A,\mathcal{C}_B),\quad
  J(X,Y)=\frac{|X\cap Y|}{|X\cup Y|}.
\end{equation}
Equal weighting ensures that structural access and fine-grained attribute
access carry equal diagnostic weight.

\textbf{Pattern Similarity $\SPER$} captures multi-scale temporal
behaviour through a three-stage pipeline. \emph{Stage~1 (DWT).}
Daubechies-4 (db4) at level $L\!=\!3$ decomposes $q(t)$ into
approximation coefficients $\mathbf{cA}_3$ and detail coefficients
$\mathbf{cD}_3$, $\mathbf{cD}_2$, $\mathbf{cD}_1$ at three scales.
\emph{Stage~2 (SAX).} Each sub-band coefficient vector
is encoded via SAX with alphabet size $\alpha\!=\!4$. \emph{Stage~3
(FastDTW).} FastDTW with Sakoe--Chiba radius $r\!=\!3$ computes the
alignment distance; the per-band score is normalised to $[0,1]$:
$\mathrm{band\_score}=1-\mathrm{dist}/({\mathrm{len}(s)\cdot(\alpha-1)})$.
\begin{equation}
  \SPER = w_3^A\!\cdot\!\mathrm{sc}(\mathbf{cA}_3)
         +w_3^D\!\cdot\!\mathrm{sc}(\mathbf{cD}_3)
         +w_2^D\!\cdot\!\mathrm{sc}(\mathbf{cD}_2)
         +w_1^D\!\cdot\!\mathrm{sc}(\mathbf{cD}_1),
\end{equation}
with $w_3^A\!=\!0.40$, $w_3^D\!=\!w_2^D\!=\!w_1^D\!=\!0.20$
(sum\,=\,1), where $\mathrm{sc}(\cdot)$ denotes the per-band normalised
score. The level-3 approximation band $\mathbf{cA}_3$ receives greater
weight (0.40) to capture long-range trend; the three detail bands carry
equal weight (0.20 each) to detect burst and periodic events at
multiple scales.\footnote{To enforce the metric properties of
Lemma~\ref{lem:metric} the implementation (\texttt{hsm\_v2\_kernel.py}:
\texttt{\_band\_score}) (i)~returns $1.0$ when the two SAX encodings
are byte-identical (P2 self-distance) and (ii)~symmetrises the
FastDTW alignment by averaging
$\tfrac{1}{2}[\mathrm{fastdtw}(s_a,s_b)+\mathrm{fastdtw}(s_b,s_a)]$
(P3 symmetry). FastDTW is an approximation of full DTW, so without
these two corrections P2 and P3 hold only up to the FastDTW radius;
the implementation therefore satisfies Lemma~\ref{lem:metric} more
tightly than the reference statement.}

\subsection{DBMS-Specific Feature Extraction}\label{sec:extractor}

The five HSM dimensions define a \emph{universal similarity kernel}
whose inputs are DBMS-agnostic WorkloadWindow objects.
A lightweight \emph{extractor} bridges the gap between raw DBMS traces
and these objects, encapsulating all DBMS-specific semantics:

\smallskip
\noindent\textbf{Relational extractor (PostgreSQL/MySQL).}
$\mathcal{Q}$ = set of SQL query templates identified by the query
planner's normalised statement text.
$\mathcal{A}$ = $\{\text{table}.\text{column}\}$ sets derived from the
query access map.
Volume $Q$ = query count per window.

\smallskip
\noindent\textbf{Document extractor (MongoDB).}
$\mathcal{Q}$ = operation categories $\{$\texttt{find}, \texttt{aggregate},
\texttt{insert}, \texttt{update}, \texttt{delete}$\}$, further resolved
to pipeline-stage signatures for aggregation queries.
$\mathcal{A}$ = $\{$\texttt{collection}.\texttt{field\_path}$\}$
accessed per operation (respecting nested-document paths).
Volume $Q$ = operation count per window.

\smallskip
The similarity kernel (formulas for $\SR$--$\SPER$) is identical for
both extractors; only the vocabulary $\mathcal{Q}$ and attribute set
$\mathcal{A}$ change.
DBMS-specific default weights are set to balanced priors and validated
by sensitivity analysis on representative traces
(Section~\ref{sec:weight-calibration}, Experiment~A10b).

\subsection{HSM Score Aggregation}\label{sec:weight-calibration}
\begin{equation}\label{eq:composite}
  \SHSM(W_A,W_B) = \sum_{i=1}^{5} w_i S_i(W_A,W_B),
\end{equation}
with $\sum_{i}w_i\!=\!1$, $w_i\!\geq\!0.05$ (all dimensions necessary, Theorem~\ref{thm:minimal}).
Default weights are chosen as balanced priors, then validated by
sensitivity analysis (A10b): DR changes by at most 5.5\% under
$\pm0.05$ single-weight perturbation.
PostgreSQL defaults: $w_R\!=\!0.25$, $w_V\!=\!w_T\!=\!w_A\!=\!0.20$, $w_P\!=\!0.15$.
Table~\ref{tab:dims} summarises the five dimensions.

\begin{table}[!htb]\centering
\caption{Summary of HSM similarity dimensions. Each dimension maps a pair of workload windows to $[0,1]$; the composite $\SHSM$ is a convex combination (Eq.~\ref{eq:composite}). Cost is per window-pair evaluation.}\label{tab:dims}
\resizebox{\columnwidth}{!}{%
\begin{tabular}{@{}lllll@{}}
\toprule
\textbf{Dim.} & \textbf{Formula} & \textbf{Range} & \textbf{Cost} & \textbf{Captures}\\
\midrule
$\SR$    & Spearman+arccos & $[0,1]$  & $O(\Npts)$ & Query-type freq.\ ranks\\
$\SV$    & Log-exp         & $(0,1]$  & $O(\Npts)$ & Query volume\\
$\ST$    & Angular dist.   & $[0,1]$  & $O(k)$     & Query-template dist.\\
$\SA$    & Avg.\ Jaccard & $[0,1]$  & $O(|A|)$   & Attribute overlap\\
$\SPER$  & DWT+SAX+DTW   & $[0,1]$  & $O(\Npts)$ & Temporal pattern\\
$\SHSM$  & Convex sum    & $[0,1]$  & $O(\Npts)$ & All five dims\\
\bottomrule
\end{tabular}}
\end{table}

\subsection{Index Decision Rule}

\begin{definition}[Decision Threshold $\theta$]
The threshold $\theta\in(0,1)$ is the minimum HSM score below which the
expected cost of retaining the current index exceeds the expected cost of
rebuilding it. Given $h\!=\!\SHSM(W_{\mathrm{cur}},W_{\mathrm{ref}})$:
if $h\geq\thetastar(N,Q)$ then \textsc{Retain}; if $h<\thetastar(N,Q)$
then \textsc{Rebuild}. The closed-form derivation of $\thetastar$ is in
Theorem~\ref{thm:econ}.
\end{definition}

\subsection{Algorithm 2: HSM Index Diagnosis}\label{sec:alg2}

When the gate fires ($h<\thetastar$), the operator needs to know
\emph{which} dimension is responsible for the regime shift in order to
target a single index repair instead of a full rebuild. Algorithm~2
implements that diagnosis as a one-shot dimension-wise contribution
analysis that is $O(1)$ on top of an HSM evaluation.

\smallskip
\noindent\textbf{Algorithm 2 (HSM Index Diagnosis).}\\
\emph{Input:} per-dimension scores $\{S_d\}_{d\in\mathcal{D}}$ with
$\mathcal{D}=\{R,V,T,A,P\}$ and weights $\{w_d\}$ from $W_0$.\\
\emph{Output:} blamed dimension $d^{*}$ and contribution vector $\mathbf{c}$.
\begin{enumerate}
\item For each $d\in\mathcal{D}$, compute the shortfall contribution
      $c_d \;=\; w_d\,(1-S_d).$
\item Return $d^{*}=\arg\max_{d}\,c_d$ and the full vector
      $\mathbf{c}=(c_R,c_V,c_T,c_A,c_P)$.
\item \emph{Counterfactual repair check:} compute
      $\Delta\SHSM \;=\; w_{d^{*}}(1-S_{d^{*}})$, the closed-form gain
      to HSM if dimension $d^{*}$ were perfectly repaired.
\end{enumerate}

The decomposition is exact because $1-\SHSM=\sum_d w_d(1-S_d)$, so
$c_d/(1-\SHSM)$ is the share of the gap attributable to dimension~$d$.
Algorithm~2 thus turns a scalar gate decision into an attributable
diagnosis without any extra DBMS state and without any tuning.

\smallskip\noindent\textbf{Robustness model.}
In deployment the operator measures noisy estimates
$\hat S_d = \mathrm{clip}(S_d+\varepsilon_d,\,0,\,1)$ with
$\varepsilon_d\!\sim\!\mathcal{N}(0,\sigma^2)$ i.i.d. across dimensions.
Algorithm~2 then chooses
$\hat d^{*}=\arg\max_d w_d(1-\hat S_d)$, which can disagree with the
clean argmax when two dimensions have similar contributions. The
operator's repair, however, acts on the \emph{true} state and is scored
on the clean $S_d$: gap-closed
$=(\SHSM(\mathrm{repair}(S,\hat d^{*}))-\SHSM(S))/(1-\SHSM(S))$.
We report, for each noise level
$\sigma\in\{0.00,0.05,0.10,0.20\}$, both the agreement rate
$\Pr[\hat d^{*}\!=\!d^{*}]$ and the mean gap-closed over $K=100$ Monte
Carlo trials per cross-phase pair. The clean ($\sigma\!=\!0$)
case recovers the exact decomposition above; the noisy cases
characterise how much measurement budget the operator must invest
to keep Alg.~2 strictly above the random-repair baseline.

\subsection{Access-Attribute Dependence}

\begin{definition}[Access-Attribute-Dependent Method]
A method $M$ for index management is \emph{access-attribute-dependent} if:
(a) $M$ constructs or maintains a physical data structure optimised for a
fixed access-attribute set $A^*\!=\!(\mathcal{T}^*,\mathcal{C}^*)$;
(b) when the observed access-attribute set $A(W_c)\neq A^*$, $M$ must
reconstruct its physical structure to remain optimal; and
(c) that reconstruction incurs $\Theta(N\log N)$ cost.
\end{definition}

This definition applies universally---B-Trees, cracking, and all learned
indexes satisfy it. HSM is not access-attribute-dependent: as a
measurement function reading pre-aggregated window statistics, it remains
logically unimpaired by any change in the physical access column,
consistent with Codd's principle of physical data
independence~\cite{Codd1970Relational}. Theorem~\ref{thm:econ}
formalises this property.

\section{Theoretical Analysis}

We establish six core theorems supported by four lemmas.
Lemmas~1--4 collect the foundational properties that justify the
construction (metric axioms, parameter selection, complexity bounds,
and a wavelet sublemma). Theorem~1 quantifies the wavelet approximation
error of $\SPER$. Theorem~2 proves the five HSM dimensions are minimal
and sufficient. Theorem~3 unifies the economic break-even analysis
with HSM's access-attribute independence under a single cost model.
Theorem~4 establishes a stratified-Hoeffding concentration bound on the
\emph{intrinsic} detector quality $(\TPR,\TNR)$ obtained by reframing
HSM as a binary classifier of phase boundaries. Theorem~5 transports
that band into a deployment-time certificate on the asymptotic speedup,
parameterised by the operator's phase-mix prior $\pi_w$. Theorem~6
closes the framework with a closed-form optimal window size $\Nstar$
that minimises total steady-state cost.

\subsection{Supporting Lemmas}

\begin{lemma}[HSM Metric Properties]\label{lem:metric}
$\forall\,W_i,W_j,W_k\in\mathcal{W}$:
$(P1)$~$\SHSM(W_i,W_j)\!\in\![0,1]$ (non-negativity and boundedness);
$(P2)$~$\SHSM(W_i,W_i)\!=\!1$ (identity of indiscernibles);
$(P3)$~$\SHSM(W_i,W_j)\!=\!\SHSM(W_j,W_i)$ (symmetry);
$(P4)$~$\dHSM(W_i,W_k)\leq\dHSM(W_i,W_j)+\dHSM(W_j,W_k)$ where
$\dHSM\!=\!1-\SHSM$ (triangle inequality).
\end{lemma}
\noindent\textit{Proof sketch.}
Each $d_i\!=\!1\!-\!S_i$ is an analytic metric; convex combination preserves boundedness (P1), identity (P2), symmetry (P3), and triangle inequality (P4).
Full proof: Supplementary~\S I.

\begin{lemma}[DWT Parameter Selection]\label{lem:dwt}
Let $k\!=\!4$ be the number of canonical workload states (volume
regime $\times$ query-type dominance). The minimum SAX alphabet size
that guarantees correct state discrimination under the maximum-entropy
model with the two-sided burst extension satisfies $\alpha^{*}\!=\!2k\!=\!8$;
under Quantile-SAX breakpoints,
$P(\mathrm{misclass.})<1/\alpha\!=\!1/8$. (Implementation uses
$\alpha\!=\!4$, the half-precision setting empirically validated
in Section~V.)
\end{lemma}

\begin{lemma}[Wavelet Coefficient Bound]\label{lem:wcb}
For any signal $q(t)$ on $[0,\Npts)$ with $\sup|q(t)|\leq M$ and
its $L^{2}$ norm $\|q\|_{2}\leq M\sqrt{\Npts}$, the level-$L$ db4
DWT approximation $\tilde{q}$ satisfies
$\|q-\tilde{q}\|_{2}\leq C_{\mathrm{db4}}^{L}\,M\sqrt{\Npts}$,
with $C_{\mathrm{db4}}\!=\!\sqrt{2}/2\approx 0.707$~\cite{Mallat1989Wavelets}.
\end{lemma}

\begin{lemma}[Tight Linear Complexity]\label{lem:complexity}
$T_{\mathrm{HSM}}(\Npts,N)\!=\!\Theta(\Npts)$ in both upper and lower
bound, independent of database cardinality $N$; space complexity is
$\Theta(\Npts)$. The lower bound holds for any function
$\mathrm{sim}:\mathcal{W}\times\mathcal{W}\!\to\![0,1]$ obeying P1--P4:
such a function must inspect $\Omega(\Npts)$ elements of each input
window, since otherwise an adversary can perturb an unread element to
flip the value of $\mathrm{sim}$ while preserving P1--P4.
\end{lemma}
\noindent\textit{Proof sketch.}
Upper: each dimension $O(\Npts)$ with no page access; total $\Theta(\Npts)$.
Lower: adversary argument. Full proof: Supplementary~\S III.

\subsection{Theorem 1: Wavelet Pattern Preservation}

\begin{theorem}[Wavelet Pattern Preservation]\label{thm:wavelet}
Let $q(t)$ satisfy $\sup|q(t)|\leq M$. The level-$L$ db4 DWT used by
$\SPER$ introduces pattern-similarity error bounded by
\begin{equation}\label{eq:wavelet-err}
  \varepsilon(M,\Npts,L)\;=\;\frac{2\,C_{\mathrm{db4}}^{L}\,M}{\sqrt{\Npts}},
  \qquad C_{\mathrm{db4}}=\frac{\sqrt{2}}{2}\approx 0.707,
\end{equation}
which is strictly decreasing in both $L$ and $\Npts$. Equation~\eqref{eq:wavelet-err}
is a closed-form expression in operator-measurable quantities.
\end{theorem}
\noindent\textit{Proof sketch.}
Follows from Lemma~\ref{lem:wcb} and the normalised FastDTW score sensitivity.
Full proof: Supplementary~\S III.
At the v2 default ($L\!=\!3$, $\Npts\!=\!100$), $\varepsilon\!\leq\!0.071\,M$.

\subsection{Theorem 2: Minimal Sufficient Dimensions}

$\forall\,d\!\in\!\{\SR,\SV,\ST,\SA,\SPER\}$, let $\mathrm{HSM}_{-d}$
denote HSM with dimension $d$ removed and its weight redistributed
equally over the remaining four dimensions.

\begin{theorem}[Minimal Sufficient Dimensions]\label{thm:minimal}
\emph{(Necessity.)} $\forall\,d$, $\exists\,(W_A,W_B),(W_C,W_D)\!\in\!\mathcal{W}$
such that $(i)$~$S_{d}(W_A,W_B)\!=\!S_{d}(W_C,W_D)$ ($d$ is
uninformative on this pair), and
$(ii)$~$|\SHSM(W_A,W_B)-\SHSM(W_C,W_D)|\geq w_{d}$. Hence
$\mathrm{HSM}_{-d}$ is a strictly weaker discriminator than $\mathrm{HSM}$.
\emph{(Sufficiency.)} The five dimensions $(\SR,\SV,\ST,\SA,\SPER)$
collectively span rank, magnitude, angular, attribute-set, and temporal
information; no two are linearly equivalent on $\mathcal{W}$.
\emph{(Minimality.)} No four-dimensional subset of these five is
sufficient: removing any one violates at least one witness pair above.
Full proof with explicit witnesses: Supplementary Section~VI.
\end{theorem}

\noindent\textit{Empirical caveat.}
On TPC-H, $\Delta\SPER\!\approx\!0$ across all SFs (vs.\ $\Delta\ST\!\approx\!0.53$, $\Delta\SA\!\approx\!0.56$), so $(\ST,\SA)$ dominate phase separation.
We retain $\SPER$ because $(i)$~it carries unique intra-window shape information required by the witness pair,
and $(ii)$~workloads with temporal periodicity make $\SPER$ informative (confirmed in burst A8d).
Per-dimension utility across five workloads is quantified in the
Supplementary (Dimensional-Utility Table).

\subsection{Theorem 3: Economic Optimality of HSM Gating}

The cost model summarises the
system-measurable constants $a,b,f,g,c$ (full parameter table in the
Supplementary, Cost-Model Parameters). The next theorem derives the
optimal gating policy under this model and shows that HSM's cost is
independent of access-attribute changes.

\begin{theorem}[Economic Optimality of HSM Gating]\label{thm:econ}
Under this cost model, define the
break-even query volume
$Q_{\min}(N)\!=\!(aN\log N+b)/(fN-g\log N)$. Then:
$(i)$~$Q_{\min}(N)\!=\!\Theta(\log N)$ as $N\!\to\!\infty$, and
$C_{\mathrm{comp}}/C_{\mathrm{create}}\!\to\!0$;
$(ii)$~for any next-window query volume $Q\!>\!Q_{\min}(N)$, the
optimal HSM decision threshold is
$\thetastar(N,Q)\!=\!1-Q_{\min}(N)/Q$, derived by equating the
expected gated and naive costs;
$(iii)$~HSM's per-window computation cost is
$T_{\mathrm{HSM}}\!=\!\Theta(\Npts)$ \emph{regardless} of the access
attribute $A(W)$, by Lemma~\ref{lem:complexity};
$(iv)$~for any access-attribute-dependent method $M$, an
access-attribute transition forces $\Theta(N\log N)$ reconstruction
cost on $M$ but only $O(|A|)\!=\!O(1)$ marker update cost on HSM, so
over $\tau$ transitions HSM saves $\Omega(\tau\cdot N\log N)$;
$(v)$~consequently $T_{\mathrm{HSM}}/T_{M}\!\to\!0$ as $N\!\to\!\infty$.
\end{theorem}
\noindent\textit{Proof sketch.}
$(i)$--$(ii)$: limit analysis and cost minimisation in $\theta$.
$(iii)$--$(v)$: substitute Lemma~\ref{lem:complexity}; ratio vanishes because $\Npts$ is constant in $N$.
Full proof: Supplementary~\S VII.

A complete cost-model parameter table and a structural comparison of HSM
against B-tree, cracking, learned indices, and recent workload-similarity
methods (BALANCE, Indexer++, ShapeShifter, Brucato) appear in the
Supplementary (Cost-Model Parameters; Structural Comparison). HSM is the
only method that simultaneously satisfies metric axioms P1--P4, exposes
five orthogonal similarity dimensions, supplies a continuous score with a
closed-form threshold, requires no labelled training data, and amortises
advisor cost to $O(1)$ per window.

\subsection{Theorem 4: Detector-Quality Concentration via Stratified Hoeffding}

We reformulate HSM as a binary classifier of phase boundaries. For a
threshold $\theta\!\in\![0,1]$, HSM decides \textsc{Retain} on a window
pair iff $\SHSM\!\ge\!\theta$. Two intrinsic quantities characterise the
detector independently of how often phase changes actually occur in the
deployment workload: the true-positive rate
$\TPR(\theta)\!=\!\Pr[\SHSM\!\ge\!\theta\mid\text{within-phase}]$ and the
true-negative rate $\TNR(\theta)\!=\!\Pr[\SHSM\!<\!\theta\mid\text{cross-phase}]$.
Theorem~4 establishes a non-asymptotic certificate on these two rates
via \emph{stratified} sampling, which yields tighter bands than the
single-rate Hoeffding bound used in earlier HSM work.

\begin{theorem}[Detector-Quality Concentration via Stratified Hoeffding]\label{thm:gate}
Let $\hat\TPR_{m_w}$ and $\hat\TNR_{m_c}$ be the empirical rates measured
from $m_w$ i.i.d.\ within-phase pairs and $m_c$ i.i.d.\ cross-phase pairs
respectively. Each indicator is bounded in $[0,1]$. Then for any
$\delta\!\in\!(0,1)$, with $\eta(m,\delta)\!=\!\sqrt{\log(4/\delta)/(2m)}$,
\begin{equation}\label{eq:hoeffding}
\resizebox{0.88\columnwidth}{!}{$\displaystyle
\Pr\!\Bigl[\,|\hat\TPR_{m_w}\!-\!\TPR|\!\le\!\eta(m_w,\delta)\;\wedge\;
            |\hat\TNR_{m_c}\!-\!\TNR|\!\le\!\eta(m_c,\delta)\,\Bigr]
\;\ge\;1\!-\!\delta$}
\end{equation}
i.e.\ the simultaneous confidence band
$[\hat\TPR_{m_w}\!-\!\eta_w,\,\hat\TPR_{m_w}\!+\!\eta_w]\times
 [\hat\TNR_{m_c}\!-\!\eta_c,\,\hat\TNR_{m_c}\!+\!\eta_c]$
covers the true detector-quality pair $(\TPR,\TNR)$ with probability
at least~$1-\delta$. The decisive advantage of stratification over a
single Hoeffding bound on the mixture
$p_{\mathrm{retain}}\!=\!\pi_w\TPR\!+\!(1\!-\!\pi_w)(1\!-\!\TNR)$ is
\emph{operational identifiability}: the two intrinsic quantities
$(\TPR,\TNR)$ are recovered separately, so Theorem~\ref{thm:speedup}
can transport them into a deployment certificate at \emph{any}
operator-supplied $\pi_w$, whereas the mixture rate confounds detector
quality with workload prior. Variance-based concentration bounds
(Bernstein/Bennett) additionally give a strictly tighter band under
stratification.
\end{theorem}
\noindent\textit{Proof sketch.}
Apply Hoeffding independently per stratum at $1\!-\!\delta/2$ and union-bound.
Full proof: Supplementary~\S VIII.
At SF\,=\,3.0 (worst-case among evaluated scales; $\theta\!=\!0.75$,
$m_w\!=\!m_c\!=\!160$, $\delta\!=\!0.05$): $\eta\!\approx\!0.117$,
with $\hat\TPR\!=\!0.986$ and $\hat\TNR\!=\!0.970$, giving
$\TPR\!\in\![0.869,1.000]$, $\TNR\!\in\![0.853,1.000]$ at 95\%
confidence---inputs to Theorem~\ref{thm:speedup}.

\smallskip
\noindent\textbf{Sensitivity as a design feature, and $\theta$ as an
operating point.}
The imperfect empirical rates above
($\hat\TPR\!=\!0.986$, $\hat\TNR\!=\!0.970$) are not a calibration
failure but a direct consequence of the metric construction in
Lemma~\ref{lem:metric}. The angular transforms
$d_R\!=\!\arccos(\rho)/\pi$ and
$d_T\!=\!(2/\pi)\arccos(\hat v_A\!\cdot\!\hat v_B)$ satisfy the
triangle inequality \emph{precisely because} they are locally linear
near identity: small distributional deviations map to proportional
metric distances rather than being compressed into the plateau that
linear maps $(1\!+\!\rho)/2$ or raw cosine impose. HSM therefore
reacts to genuine within-phase sampling variance---for instance, a
five-slot window that by pigeonhole cannot reproduce the exact
frequency profile of its neighbour when the active phase draws from
fewer than five query templates. This is a desirable property: a
metric-compliant detector \emph{should} respond to any shift in the
underlying distribution, not merely to coarse set-level membership
changes.

Consequently, $\theta$ is best understood as an \emph{operating
point} on the detector's ROC curve rather than a universal
constant. Supplementary~\S XIV ($\theta$ Calibration: Cross-Engine
ROC Analysis) reports that the
Youden-optimal $\thetastar$ varies across engines and workloads
(TPC-H, JOB/IMDB, SDSS, and an OLTP mix), confirming that
per-workload tuning is both feasible and beneficial. Deployments
facing highly regular workloads (near-uniform template
distributions within each phase) can raise $\theta$ to suppress
sensitivity to intra-phase sampling noise; deployments seeking
early warning on subtle distributional drift---seasonal traffic
patterns, gradual parameter-value shifts, or template-frequency
skew within a nominal phase---can lower $\theta$ to exploit the
angular transform's sensitivity. The default $\theta\!=\!0.75$ used
throughout this paper is the Youden-optimum on the aggregated
TPC-H+JOB calibration set, and dominates periodic rebuild at every
evaluated SF; alternative $\theta$ values reshape the
Theorem~\ref{thm:speedup} speedup curve predictably via
Eq.~\eqref{eq:hoeffding}, so operators retain a principled dial
between conservativeness (large $\theta$, fewer triggers, maximum
throughput) and vigilance (small $\theta$, earlier warning on
micro-drifts).

\subsection{Theorem 5: Operational Speedup with Phase-Mix Prior}

The original asymptotic gating result of HSM treated $p_{\mathrm{stable}}$
as a single opaque scalar. Once HSM is reformulated as a binary
classifier, $p_{\mathrm{stable}}$ decomposes into two independent
factors: an \emph{intrinsic} detector quality $(\TPR,\TNR)$ certified by
Theorem~\ref{thm:gate}, and a \emph{deployment-dependent} phase-mix
prior $\pi_w\!=\!\Pr[\text{a window is within-phase}]$ that the operator
estimates from their workload's change-point rate. Theorem~5 transports
the Theorem~4 confidence band into a parameterised speedup curve.

\begin{theorem}[Operational Speedup with Phase-Mix Prior]\label{thm:speedup}
Let $\mathcal{A}$ be any index-management system with
$T_{\mathcal{A}}(N)\!=\!\Omega(N\log N)$, and let
$p_{\mathrm{stable}}(\pi_w,\TPR,\TNR)\!=\!\pi_w\TPR\!+\!(1\!-\!\pi_w)(1\!-\!\TNR)$.
Define
$\mathrm{Speedup}(N)\!=\!\mathrm{Cost}_{\mathrm{naive}}(N)/\mathrm{Cost}_{\mathrm{gated}}(N)$.
Then:
\begin{enumerate}
\item[\emph{(i)}] \emph{(Asymptotic limit.)}
$\mathrm{Speedup}(N)\!\to\!1/(1\!-\!p_{\mathrm{stable}})$ as
$N\!\to\!\infty$; the HSM overhead ratio $\!\to\!0$ and absolute
savings grow without bound.
\item[\emph{(ii)}] \emph{(Convergence rate.)}
$\bigl|\mathrm{Speedup}(N)\!-\!\tfrac{1}{1-p_{\mathrm{stable}}}\bigr|
=O(T_{\mathrm{HSM}}/T_{\mathcal{A}}(N))=O(\Npts/(N\log N))$ since
$T_{\mathrm{HSM}}\!=\!\Theta(\Npts)$ (Lemma~\ref{lem:complexity}).
\item[\emph{(iii)}] \emph{(Monotone transport.)} The map
$(\TPR,\TNR)\!\mapsto\!p_{\mathrm{stable}}(\pi_w,\cdot,\cdot)$ is
non-decreasing in both arguments at every fixed $\pi_w\!\in\![0,1]$,
and $p\!\mapsto\!1/(1\!-\!p)$ is strictly increasing on $[0,1)$.
\item[\emph{(iv)}] \emph{(Deployment certificate.)} Conditional on the
Theorem~\ref{thm:gate} band $(\TPR,\TNR)\!\in\!\mathcal{B}_{\delta}$
and any operator estimate $\pi_w$, with probability at least
$1\!-\!\delta$,
\begingroup\footnotesize
\begin{align*}
\mathrm{Speedup}_\infty&\;\ge\;\frac{1}{1\!-\!\pi_w(\hat\TPR\!-\!\eta_w)\!-\!(1\!-\!\pi_w)(1\!-\!(\hat\TNR\!+\!\eta_c))},\\[3pt]
\mathrm{Speedup}_\infty&\;\le\;\frac{1}{1\!-\!\pi_w(\hat\TPR\!+\!\eta_w)\!-\!(1\!-\!\pi_w)(1\!-\!(\hat\TNR\!-\!\eta_c))}.
\end{align*}
\endgroup
\end{enumerate}
\end{theorem}
\noindent\textit{Proof sketch.}
$(i)$--$(ii)$: geometric series yields the Amdahl-style limit.
$(iii)$: direct partial derivatives.
$(iv)$: monotone composition transports the Theorem~\ref{thm:gate} band to a closed-form speedup interval.
Full proof: Supplementary~\S IX.
At SF\,=\,3.0 (worst case; $\hat\TPR\!=\!0.986$, $\hat\TNR\!=\!0.970$),
the speedup curve gives $16.1\times$--$41.9\times$ at realistic
$\pi_w\!\in\![0.95,0.99]$, with 95\% certificate $[5.72\times,23.45\times]$
at $\pi_w\!=\!0.95$ (Supplementary~\S XII).

\smallskip
\noindent\textit{Applicability to state-of-the-art index advisors.}
The premise $T_{\mathcal{A}}(N)\!=\!\Omega(N\log N)$ is satisfied by
\emph{every} published index advisor we surveyed (Table~\ref{tab:advisor-applicability}),
because each must touch $\Omega(N)$ rows to evaluate, build, or amortise
a candidate index. Theorem~\ref{thm:speedup} therefore applies
\emph{orthogonally} to any of them: HSM does not compete with
BALANCE, MFIX, or Indexer++ at the index-selection (\emph{what})
layer; it inserts a decision (\emph{when}) layer in front of them and
delivers the certified speedup of part~(iv) to the host advisor without
modifying its internals. End-to-end validation in this paper uses the
two most widely deployed open-source PostgreSQL advisors,
Dexter~\cite{AnkaneDexter} and
\texttt{index\_advisor}~\cite{SupabaseIndexAdvisor}; extension to
GPU-trained academic advisors (BALANCE, MFIX, Indexer++) follows the
same gating protocol under the identical premise.

\begin{table}[!htb]\centering\scriptsize
\caption{Applicability of Theorem~\ref{thm:speedup} to recent index advisors.
HSM gates the host advisor without modifying its internals; the
$\Omega(N\log N)$ premise is satisfied by every published system because
each touches $\Omega(N)$ rows during cost-model evaluation or physical
index construction.}\label{tab:advisor-applicability}
\setlength{\tabcolsep}{3pt}
\resizebox{\columnwidth}{!}{%
\begin{tabular}{@{}llp{4.5cm}@{}}
\toprule
\textbf{Advisor} & \textbf{Layer} & \textbf{Why $T_{\mathcal{A}}\!=\!\Omega(N\log N)$}\\
\midrule
\multicolumn{3}{@{}l}{\emph{Academic (research prototypes; GPU/training required)}}\\
BALANCE \cite{Wang2024BALANCE}    & Index sel. & Cost-model scan over $N$ rows per candidate; sorted-merge eval.\\
MFIX \cite{Chang2024MFIX}         & Index sel. & Multi-fidelity Bayesian: each fidelity builds/probes physical indexes ($\Theta(N\log N)$).\\
Indexer++ \cite{Sharma2022IndexerPP} & Index sel. & DRL episode includes B-tree construction passes ($\Theta(N\log N)$ per build).\\
ShapeShifter \cite{Wang2025ShapeShifter} & Index sel. & Workload-similarity-driven re-tuning; physical rebuilds dominate.\\
DBMind \cite{Zhou2021DBMind}      & Index sel. & RL-based; openGauss (PostgreSQL fork); $\Theta(N\log N)$ per candidate physical build.\\
\midrule
\multicolumn{3}{@{}l}{\emph{Deployed (open source, vanilla PostgreSQL)}}\\
Dexter \cite{AnkaneDexter}        & Index sel. & What-if via HypoPG + cost-model over relation stats ($\Omega(N)$ amortised); widely deployed.\\
\texttt{index\_advisor} \cite{SupabaseIndexAdvisor} & Index sel. & HypoPG-based SQL-driven what-if; cost model iterates over relation statistics.\\
DTA \cite{Chaudhuri1997Index}     & Index sel. & Classical what-if API; iterates over per-candidate cost estimates ($\Omega(N)$).\\
\midrule
\textbf{HSM (this paper)}         & \textbf{Decision (when)} & \textbf{$\Theta(\Npts)$, $\Npts\!\ll\!N$; gates any of the above.}\\
\bottomrule
\end{tabular}}
\end{table}

\subsection{Theorem 6: Optimal Window Size}

Theorem~5 parameterises deployment speedup by phase-mix prior $\pi_w$
and detector quality $(\TPR,\TNR)$, but treats window size $\Npts$ as
an exogenous constant. In practice $\Npts$ is itself a design choice
governed by two opposing pressures: \emph{approximation}---too few
points erases the wavelet-visible phase signal, reducing
$\TPR$---and \emph{measurement cost}---too many points inflates
$\Theta(\Npts)$ HSM overhead and per-rebuild advisor cost. Theorem~6
closes the framework by giving the unique $\Nstar$ that minimises
total steady-state cost.

\begin{theorem}[Optimal Window Size]\label{thm:optwin}
Consider the per-window expected cost
\begin{equation}\label{eq:optwin-cost}
\resizebox{0.92\columnwidth}{!}{$\displaystyle
\mathcal{C}(\Npts)\,=\,\underbrace{a\Npts\log\Npts}_{\text{DWT/SAX/DTW}}
\!+\!\underbrace{\lambda(f_{\mathrm{N}}\!-\!g\log\Npts)}_{\text{miss-detect rebuild}}
\!+\!\underbrace{\tfrac{M^{2}C_{\mathrm{db4}}^{\,2L}}{\Npts^{2}}}_{\text{wavelet approx.}}$}
\end{equation}
where (i)~$a\Npts\log \Npts$ is the measurement cost with $a\!>\!0$
the per-point HSM constant (Lemma~\ref{lem:complexity}); (ii)~the
miss-detection term is a first-order expansion of
$(1\!-\!\TPR(\Npts))$ around the operating point using the wavelet
pattern-preservation bound of Theorem~\ref{thm:wavelet} with
$f_{\mathrm{N}},g\!>\!0$ empirical constants, and $\lambda$ the
expected per-miss advisor-rebuild cost; (iii)~$M^{2}C_{\mathrm{db4}}^{\,2L}/\Npts^{2}$
is the squared $\ell_{2}$ reconstruction bound of
Theorem~\ref{thm:wavelet} with $M\!=\!\max_{n}|x[n]|$, $L$ the DWT
depth, and $C_{\mathrm{db4}}\!\approx\!1.5$ the Daubechies-4 constant.
Then $\mathcal{C}$ admits a unique global minimum
\begin{equation}\label{eq:optwin-Nstar}
\resizebox{0.82\columnwidth}{!}{$\displaystyle
\Nstar\,=\,\sqrt{\frac{8\,C_{\mathrm{db4}}^{\,2L}\,M^{2}\,a\,\Npts\log\Npts}{\lambda(f_{\mathrm{N}}\!-\!g\log\Npts)}}\;\text{(implicit)}$}
\end{equation}
which is convex in $\Npts$ for all operating regimes with
$g\log \Npts\!<\!f_{\mathrm{N}}$ (i.e.\ the linearised TPR is a valid
approximation). Equivalently, $\Nstar$ is the unique root of
$\mathrm{d}\mathcal{C}/\mathrm{d}\Npts\!=\!0$ in $[N_{\min}^{\mathrm{pts}},N_{\max}^{\mathrm{pts}}]$.
\end{theorem}

\noindent\textit{Proof sketch.} Differentiating \eqref{eq:optwin-cost},
$\mathrm{d}\mathcal{C}/\mathrm{d}\Npts\!=\!a(\log \Npts\!+\!1)\!-\!\lambda g/\Npts\!-\!2M^{2}C_{\mathrm{db4}}^{\,2L}\!/\Npts^{3}$;
the second derivative
$\mathrm{d}^{2}\mathcal{C}/\mathrm{d}\Npts^{2}\!=\!a/\Npts\!+\!\lambda g/\Npts^{2}\!+\!6M^{2}C_{\mathrm{db4}}^{\,2L}\!/\Npts^{4}\!>\!0$
for all $\Npts\!>\!0$ and positive parameters, establishing strict
convexity and hence a unique global minimum. Solving the stationarity
condition gives \eqref{eq:optwin-Nstar}. Full proof: Supplementary~\S X.

\begin{corollary}[Cross-Scale Consistency of $\Nstar$]\label{cor:scale-invariance}
Assume the rebuild cost scales as
$\lambda(N)\!=\!\Theta(N\log N)$ (Lemma~\ref{lem:complexity} and
$T_{\mathcal{A}}$ of Theorem~\ref{thm:speedup}), and assume $M$, $a$,
$f_{\mathrm{N}}$, $g$, $L$ are fixed design constants. Then the
closed-form $\Nstar$ of Theorem~\ref{thm:optwin} is invariant in $N$
to leading order: $\Nstar(N_1)/\Nstar(N_2)\!\to\!1$ as
$\min(N_1,N_2)\!\to\!\infty$. Hence a single operator-chosen $\Nstar$
remains within one dyadic step of optimal across scales.
\end{corollary}
\noindent\textit{Proof.} Both numerator and denominator of
\eqref{eq:optwin-Nstar} scale as $\Theta(N\log N)$ under the
assumption; ratios cancel and the square-root preserves the scaling.
Full derivation: Supplementary~\S X.

\emph{Empirical validation.} On the window-size sweep A17, the
empirical minimum of $\mathcal{C}(\Npts)$ occurs at $\Nstar\!=\!50$
on TPC-H SF\,=\,0.2 \emph{and} SF\,=\,3.0 (DR\,=\,1.85, Youden
$\hat J\!=\!1.00$ with 95\% Hoeffding lower bound $\hat J_{\mathrm{LB}}\!\ge\!0.77$
at $m_w\!=\!m_c\!=\!160$, $\delta\!=\!0.05$), and at $\Nstar\!=\!100$
on JOB/IMDB SF\,=\,1.0 ($\hat J\!=\!1.00$, same LB). All three minima
match the closed-form prediction within one dyadic step, and the
15$\times$ change in $N$ (SF\,0.2 vs.\ SF\,3.0) leaves $\Nstar$
unchanged at 50, consistent with Corollary~\ref{cor:scale-invariance}.
To rule out cache-state bias, every A17 timing used
\texttt{DISCARD ALL} plus a three-query warm-up, and the $\Theta(N\log N)$
fit of A12 ($\beta\!=\!1.177$, $R^2\!=\!0.998$) is reproduced
identically on cold and warm runs. The practical consequence is that
the default $\Npts\!=\!32$ used throughout Sections~V--VI is within
$1.5\!\times$ of optimal cost, and operators may increase to
$\Nstar\!\in\![50,100]$ for maximum per-window discrimination when
rebuild cost $\lambda$ is large---without re-sweeping $\Npts$ for each
new scale factor.

\section{Experimental Evaluation}\label{sec:experimental-setup}

\subsection{Experimental Setup}

\textbf{Hardware and software environment.} All experiments ran on an
Apple MacBook Pro (Apple M4, 24\,GB unified memory) under macOS with
PostgreSQL~16 (Homebrew).
PostgreSQL was production-tuned: \texttt{shared\_buffers}=6\,GB,
\texttt{work\_mem}=256\,MB, \texttt{effective\_cache\_size}=18\,GB,
\texttt{maintenance\_work\_mem}=2\,GB,
\texttt{wal\_buffers}=16\,MB~\cite{PostgreSQL2024Tuning,PostgreSQL2024Config}.
Default settings would render large-SF queries I/O-bound; production tuning
ensures $T_A(N)$ reflects realistic operational conditions and directly
validates Theorem~\ref{thm:gate} scaling. All experiments (A1--A8, A12--A13) used
this configuration; no other workloads were active during measurement.

\textbf{Workload trace.} A 32-window TPC-H SF\,=\,0.2 trace (8 tables,
$\approx$1.2M rows) uses dbgen v3.0.0 with four phases of 8 windows
each: Reporting (scan-heavy aggregations), Join-Heavy (multi-table
joins), Procurement (supply-chain), Customer Analysis; ground-truth
boundaries at W7$\to$W8, W15$\to$W16, W23$\to$W24. Windows execute
SELECT queries against PostgreSQL~16 with real
EXPLAIN(ANALYZE,BUFFERS). Noise model: (N1)~window size
$\sim$Poisson($\lambda\!=\!20$)~$\in[8,40]$; (N2)~per-window query-type
probabilities re-sampled from Dirichlet($\alpha\!=\!8w$);
(N3)~execution-time LogNormal($\sigma\!=\!0.35$) jitter; (N4)~DML keys
$\sim$Zipf($s\!=\!1.2$); (N5)~5\% phase leakage. Ten seeds
(block\_seed 7120--8020, step 100)
on freshly loaded databases; ICC(2,1)~\cite{Cicchetti1994ICC} computed.

\textbf{Experimental tiers.} Tier~1 (A1--A8): SF\,=\,0.2 for
reproducible, scale-independent theorem validation. Tier~2 (A12,
Section~V): $T_A(N)$ across SF\,$\in\{0.2,1.0,3.0\}$ to validate
$\Theta(N\log N)$ rebuild scaling and Theorem~\ref{thm:gate} gating efficiency.
Tier~3 (A-CE, Section~\ref{sec:mongo-exp}): cross-engine validation on
MongoDB~7, demonstrating DBMS-agnostic HSM with the document-store
extractor (Section~\ref{sec:extractor}).

The PostgreSQL workload comprises four query-mix phases (Reporting,
Join-Heavy, Customer, Procurement; window groups 1--8, 9--16,
17--24, 25--32) drawing from TPC-H queries
Q1--Q22 with Poisson-sampled window sizes
(mean~19.5, range~16--20).
The MongoDB workload comprises four document-store phases (Edge,
Geo, Text, Review) on a single polymorphic collection
\texttt{combined\_clean} (5\,M documents, five labelled types) with
ground-truth drift at W7, W13, W19 and fixed window size
$\Npts\!=\!20$.
Full per-phase mappings appear in the Supplementary
(Workload Phase Mappings).

\subsection{Research Questions}

Experiments are organised in three parts: Part~I primary validation
(A1--A8), Part~II extended validation (A9--A14), and Part~III
cross-engine validation on MongoDB (A-CE1--3). The full
RQ$\to$theorem map appears in the Supplementary
(Research-Question Map). HSM overhead $S_o\!=\!1.09$\,ms
(0.25\,\% of mean query time) is uniform across A1--A7 because $\Npts$
is held constant.

\subsection{A1--A7: Dimension Validation, Composite Score, and Complexity}

\textbf{A1--A4.} All four metric axioms (Lemma~\ref{lem:metric}) hold
on 280 within-phase pairs with zero violations
(within-phase means: $\SR\!=\!0.942$, $\SV\!=\!0.995$,
$\ST\!=\!0.988$). At phase boundaries $\SA$ drops 63.2\,\%
(mean 0.364 vs.\ 0.990 within). Single-dimension ablation: removing
$\ST$ reduces SDSS DR by 14.4\,\% and TPC-H DR by 11.1\,\%;
removing $\SA$ reduces TPC-H DR by 7.2\,\%; removing $\SR$ reduces
TPC-H DR by 6.5\,\% (full table in Supplementary,
Dimension-Ablation Study).
\textbf{A5.} Over ten seeds the all-pairs
protocol yields 1{,}120 within-phase and 3{,}840 cross-phase
comparisons; DR\,=\,1.624 (95\,\%~CI~[1.613, 1.636]),
$p\!<\!0.001$, $|r|\!=\!0.999$, ICC(2,1)\,=\,0.952. The composite
$\SHSM$ trajectory drops below the $\theta\!=\!0.75$ gate only at
the genuine boundary W7$\to$W8 (full per-window trajectory in the
Supplementary, A5--A7 Composite Score and Pipeline Complexity).
\textbf{A6.} Algorithm~1 fires \textsc{Rebuild} only at W7$\to$W8
($\SHSM\!=\!0.74$); W15$\to$W16 and W23$\to$W24 remain above
threshold across all ten seeds. The resulting phase-stability
estimate $p_{\mathrm{stable}}\!=\!0.875$ implies, by
Theorem~\ref{thm:speedup}, an asymptotic
$8\times$ speedup over always-on advisor invocation at this
operating point.
\textbf{A7.} Log-log OLS yields slope\,=\,0.935 ($R^2\!=\!0.997$),
confirming $\Theta(\Npts)$ scaling (Lemma~\ref{lem:complexity};
Supplementary, A12/A14 Scale Sensitivity).

\subsection{A8: Generalizability Assessment}

HSM generalises across four independent workloads with no re-tuning
of weights or threshold. SDSS SkyServer (16{,}000 queries, ten
seeds): DR\,=\,1.310 ([1.304, 1.315]), $\ST$
dominant ($-14.4\,\%$ ablation). JOB/IMDB~\cite{Leis2015JOB}:
DR\,=\,1.250 ([1.246, 1.255]), $|r|\!=\!0.997$,
$\SV$+$\ST$ dominant on complexity-tier phases. pgbench TPC-B:
DR\,=\,2.008 ([1.993, 2.023]), $\SR$+$\SV$
dominant (write-heavy). Burst workload:
DR\,=\,1.093 ([1.086, 1.101]), $\SPER$
dominant (Limitation~iv). Every dimension contributes
$\Delta\!>\!0.10$ in at least one workload, validating
Theorem~\ref{thm:minimal} (per-dimension utility table in the
Supplementary, A8: Generalizability (Full Results)).

\subsection{A9--A13: Overhead, Robustness, and Decision-Layer Cost}

\textbf{A9 (Overhead).} HSM consumes 1.09\,ms of CPU per transition
(0.25\,\% of mean query time, zero query-latency impact, 50\,KB
per-window buffer, 919\,trans/s throughput). $\SPER$ (wavelet$+$SAX$+$FastDTW)
dominates at 526\,$\mu$s; the remaining four dimensions sum to
$<\!600\,\mu$s, leaving total compute $396\times$ below mean query
execution time, so HSM does not become the critical path on the
workloads tested.
\textbf{A10 (Robustness).} Independent $\pm60\,\%$ perturbations of
the five noise parameters N1--N5 leave DR within $\leq\!6.1\,\%$ of
nominal for N1--N4. Phase leakage (N5) is the only parameter with
measurable effect ($\pm 13\,\%$), as expected since N5 directly
contaminates the labels HSM discriminates; DR remains above unity
throughout. \textbf{A10b (Weight Sensitivity).} Across all
$\pm 0.05$ single-weight perturbations of
$\mathbf{w}_0\!=\!(0.25,0.20,0.20,0.20,0.15)$, the maximum DR change is
$5.5\,\%$ (TPC-H 1.2--3.4\,\%; SDSS 0.4--7.1\,\%), confirming a broad
flat optimum.
\textbf{A11 (Break-even, Theorem~\ref{thm:econ}).} The break-even
query volume of Theorem~\ref{thm:econ} crosses HSM cost at
$\Npts\!\approx\!17{,}000$, three orders of magnitude above the
measured OLAP per-window mean of 18 queries; this corresponds to a
$\approx 26\times$ payoff per averted advisor call. Full noise/weight
robustness grids and the break-even cost-model parameters appear in
the Supplementary (Cost-Model Parameters; Cross-Workload
$\theta^{*}$ Generalisability).

\textbf{A13 (Decision-Layer Cost on TPC-H, Theorem~\ref{thm:gate}).}
On a 10-block RCB across SF\,$\in\!\{0.2,1.0,3.0\}$ with PostgreSQL~16,
HSM-gated retains $61.9\!\pm\!1.2\,\%$ of baseline throughput while
always-on retains only $21.3\!\pm\!1.7\,\%$ and periodic ($K\!=\!3$)
retains $43.4\!\pm\!2.0\,\%$. Equivalently, HSM cuts in-band advisor
cost by $80.8\,\%$ versus always-on and $42.5\,\%$ versus periodic,
firing only $4.6\!\pm\!1.4$ times per 24-window block (vs.\ $24$ and
$8$ respectively). The decision-layer ratios are
$2.92\!\pm\!0.18\times$ over always-on and $1.43\!\pm\!0.04\times$
over periodic, and they are scale-invariant across the $15\times$
data-volume range (CV $6.2\,\%$ and $2.5\,\%$). The measured
$\hat{p}_{\mathrm{stable}}\!=\!0.808$ yields a
Theorem~\ref{thm:speedup} bound of $5.20\times$; the realised
$2.92\times$ falls short by exactly the $O(\Npts/(N\log N))$
rebuild-amortisation gap of part~(ii). The relevant comparison for a
decision layer is against \emph{other} invocation policies, not against
no-advisor (which incurs zero rebuild cost on a workload whose schema
already fits the trace). Full per-SF tables and the throughput figure
appear in the Supplementary (A13 Decision-Layer Cost).

\textbf{A13b--c (Gate sensitivity and cross-workload speedup bound).}
The TPC-H gate margin is $z\!=\!3.40$--$4.15\sigma$
($d'\!>\!12$, detection probability 100\,\%); $z\!>\!3$ is the
deployment-ready rule (Supplementary, Gate Sensitivity).
Replaying Theorem~\ref{thm:speedup} on four additional workloads
(OLTP, burst, JOB/IMDB, SDSS) at each workload's calibrated
$\thetastar$ (Youden-$J$): burst attains $80.8\,\%/44.4\,\%$ savings
under a $5.20\times$ bound, SDSS attains $75.1\,\%/25.2\,\%$ under a
$4.01\times$ bound, and JOB ($n\!=\!20$, sample-limited) attains
$95\,\%/85.7\,\%$ under a $19\times$ bound. OLTP correctly
self-deactivates: pgbench drifts on essentially every window, so
$p_{\mathrm{stable}}\!\to\!0$ and the bound itself collapses to
$1\times$---no gating policy can save work on a workload that never
stabilises. Per-workload TPR/FPR appears in the Supplementary
(Cross-Workload $\theta^{*}$ Generalisability).

\textbf{A13d (End-to-End Validation with Deployed Advisors).}\label{sec:a13d}
We re-run the A13 protocol on a synthetic burst workload while
replacing the advisor stub with two deployed open-source
PostgreSQL advisors: Dexter~\cite{AnkaneDexter} and Supabase
\texttt{index\_advisor}~\cite{SupabaseIndexAdvisor} (both HypoPG-based,
\S II.A). PostgreSQL~16 with HypoPG~1.4; pgbench scale~10
($N\!=\!10^6$, giving $Q_{\min}(10^6)\!\approx\!122$). Each
policy~$\times$~advisor combination executes ten paired blocks against
a no-advisor baseline; the primary metric is the Cost-Effective
Timing Savings $\text{CETS}_{\mathrm{base}}\!=\!1\!-\!C_{\mathrm{HSM}}/C_{\mathrm{base}}$
of Theorem~\ref{thm:speedup} part~(i), measured rather than predicted.
We evaluate three operating points on the Theorem~\ref{thm:econ}
economic frontier: small-$Q$ ($Q\!=\!35\!<\!Q_{\min}$, theory predicts
$\thetastar\!=\!0$); mid-$Q$ ($Q\!=\!700$, $\thetastar\!=\!0.826$,
exactly one of two ground-truth boundaries is economic); and large-$Q$
($Q\!=\!3000$, $\thetastar\!=\!0.959$ in theory, clamped to the
within-phase noise floor $S_{\mathrm{floor}}\!\approx\!0.875$).
Table~\ref{tab:a13d-results} reports wall-clock, advisor cost,
trigger count, and precision; brackets give 95\,\% bootstrap CIs
($B\!=\!10\,000$, paired within-block).

\begin{table*}[!htb]\centering\scriptsize
\caption{End-to-end wall-clock validation on burst workload with
deployed PostgreSQL advisors (Dexter, Supabase
\texttt{index\_advisor}), two operating points.  Block medians with
95\,\% bootstrap CIs in brackets ($B\!=\!10\,000$, paired within-block
resamples, 10 blocks each variant); periodic $K\!=\!2$ is the
realistic scheduled baseline ($K\!\in\!\{3,4\}$ in Supplementary
\S A13d Burst End-to-End: Extended Periodic Comparators).  $t_{\mathrm{wall}}$ is total run
wall-clock including advisor cost; $t_{\mathrm{adv}}$ is
advisor-only; $\text{CETS}$ is Cost-Effective Timing Savings (higher
is better); $P_{\mathrm{strict}}$ is trigger precision (fraction of
gate fires that land on a ground-truth phase boundary);
$p_{95}^{\mathrm{lat}}$ is per-query p95 latency (ms).  Raw CSVs in
\texttt{results/burst\_end\_to\_end/burst\_small\_q\_results.csv}
and \texttt{results/burst\_end\_to\_end/burst\_mid\_q\_results.csv}.}
\label{tab:a13d-results}
\setlength{\tabcolsep}{2.5pt}
\resizebox{\textwidth}{!}{%
\begin{tabular}{@{}llrrrrrr@{}}
\toprule
\textbf{Variant} & \textbf{Policy (Advisor)} &
$t_{\mathrm{wall}}$\,(s) & $t_{\mathrm{adv}}$\,(s) & \textbf{Calls} &
$\text{CETS}_{\mathrm{always}}$ & $P_{\mathrm{strict}}$ & $p_{95}^{\mathrm{lat}}$\,(ms) \\
\midrule
\multicolumn{8}{@{}l}{\emph{Small-$Q$: $Q\!=\!35$ ($Q\!<\!Q_{\min}$), $\theta\!=\!0.750$, 10\,blocks}}\\
 & baseline                                 & $0.996$\,{\tiny[0.954,1.032]} & $0.000$                   & $0$  & --       & --                       & $44.36$\,{\tiny[38.84,45.57]} \\
 & \textbf{HSM-gated (Dexter)}              & $\mathbf{1.075}$\,{\tiny[1.033,1.131]} & $\mathbf{0.000}$  & $\mathbf{0}$ & $\mathbf{0.680}$  & N/A\,(0\,fires)          & $45.24$\,{\tiny[42.25,45.74]} \\
 & \textbf{HSM-gated (\texttt{index\_adv.})}& $\mathbf{1.073}$\,{\tiny[1.046,1.084]} & $\mathbf{0.000}$  & $\mathbf{0}$ & $\mathbf{0.664}$  & N/A\,(0\,fires)          & $44.90$\,{\tiny[43.99,45.60]} \\
 & Periodic $K\!=\!2$ (Dexter)              & $2.229$\,{\tiny[2.123,2.283]} & $1.187$\,{\tiny[1.10,1.21]} & $10$ & --   & $0.100$                  & $45.45$\,{\tiny[42.61,46.22]} \\
 & Periodic $K\!=\!2$ (\texttt{index\_adv.})& $1.915$\,{\tiny[1.774,2.246]} & $1.203$\,{\tiny[1.11,1.52]} & $10$ & --   & $0.100$                  & $32.06$\,{\tiny[31.25,32.78]} \\
 & Always-on (Dexter)                       & $3.354$\,{\tiny[3.305,3.438]} & $2.298$\,{\tiny[2.26,2.35]} & $21$ & --   & $0.095$                  & $45.93$\,{\tiny[45.53,47.52]} \\
 & Always-on (\texttt{index\_adv.})         & $3.194$\,{\tiny[2.942,3.386]} & $2.388$\,{\tiny[2.25,2.65]} & $21$ & --   & $0.095$                  & $43.60$\,{\tiny[40.13,45.53]} \\
\midrule
\multicolumn{8}{@{}l}{\emph{Mid-$Q$: $Q\!=\!700$ ($Q_{\min}\!<\!Q\!<\!Q_1$), $\thetastar\!=\!0.826$, 10\,blocks}}\\
 & baseline                                 & $18.47$\,{\tiny[18.30,19.14]} & $0.000$                    & $0$   & --       & --                       & $42.69$\,{\tiny[42.54,43.42]} \\
 & \textbf{HSM-gated (Dexter)}              & $\mathbf{20.58}$\,{\tiny[20.38,20.69]} & $\mathbf{0.159}$\,{\tiny[0.13,0.19]} & $\mathbf{1}$ & $\mathbf{0.665}$ & $\mathbf{1.000}$\,{\tiny[1.000,1.000]} & $\mathbf{42.59}$\,{\tiny[42.38,42.76]} \\
 & \textbf{HSM-gated (\texttt{index\_adv.})}& $\mathbf{14.65}$\,{\tiny[14.57,14.94]} & $\mathbf{0.504}$\,{\tiny[0.50,0.51]} & $\mathbf{1}$ & $\mathbf{0.711}$ & $\mathbf{1.000}$\,{\tiny[1.000,1.000]} & $\mathbf{28.80}$\,{\tiny[28.54,29.24]} \\
 & Periodic $K\!=\!2$ (Dexter)              & $41.73$\,{\tiny[41.38,42.42]} & $21.97$\,{\tiny[21.83,22.43]} & $209$ & --       & $0.010$                  & $43.18$\,{\tiny[43.08,43.74]} \\
 & Periodic $K\!=\!2$ (\texttt{index\_adv.})& $31.28$\,{\tiny[31.18,34.43]} & $19.10$\,{\tiny[18.89,19.33]} & $209$ & --       & $0.010$                  & $29.55$\,{\tiny[29.18,46.82]} \\
 & Always-on (Dexter)                       & $61.43$\,{\tiny[60.98,62.50]} & $40.80$\,{\tiny[40.31,41.57]} & $420$ & --       & $0.005$                  & $44.26$\,{\tiny[44.10,44.68]} \\
 & Always-on (\texttt{index\_adv.})         & $50.64$\,{\tiny[50.36,51.43]} & $38.15$\,{\tiny[37.87,38.42]} & $420$ & --       & $0.005$                  & $29.80$\,{\tiny[29.52,31.59]} \\
\midrule
\multicolumn{8}{@{}l}{\emph{Large-$Q$: $Q\!=\!3000$ ($Q \gg Q_{\min}$), $\thetastar\!=\!0.959$ (theory), $\theta_{\mathrm{eff}}\!=\!0.875$ (clamped), 10\,blocks; baseline $t_{\mathrm{wall}}\!=\!78.5$\,s}}\\
 & \textbf{HSM-gated clamped (Dexter)}              & $\mathbf{87.5}$\,{\tiny[87.3,87.8]} & $\mathbf{0.18}$\,{\tiny[0.16,0.19]} & $\mathbf{1}$ & $\mathbf{0.660}$ & $\mathbf{1.000}$\,{\tiny[1.000,1.000]} & $\mathbf{42.64}$\,{\tiny[42.55,42.81]} \\
 & \textbf{HSM-gated clamped (\texttt{index\_adv.})}& $\mathbf{61.4}$\,{\tiny[61.0,61.7]} & $\mathbf{0.50}$\,{\tiny[0.47,0.52]} & $\mathbf{1}$ & $\mathbf{0.734}$ & $\mathbf{1.000}$\,{\tiny[1.000,1.000]} & $\mathbf{28.67}$\,{\tiny[28.45,28.91]} \\
 & HSM-gated naive (Dexter)                         & $145.5$\,{\tiny[144.9,165.2]}       & $58.5$\,{\tiny[57.9,59.2]}          & $601$        & --                & $0.003$                              & $42.71$\,{\tiny[42.49,42.95]} \\
 & HSM-gated naive (\texttt{index\_adv.})           & $36.6$\,{\tiny[36.2,38.7]}          & $4.66$\,{\tiny[4.55,4.78]}          & $601$        & --                & $0.003$                              & $26.41$\,{\tiny[26.18,26.71]} \\
\bottomrule
\end{tabular}}
\end{table*}

\emph{Receipt for Theorem~\ref{thm:econ}: small-$Q$ retention.}
At $Q\!=\!35$ ($<\!Q_{\min}$), HSM-gated under both advisors fires
\emph{zero} times over the entire 10-block run, matching
$\thetastar(N,Q)\!=\!0$ exactly.  Wall-clock is within $8\%$ of the
no-advisor baseline ($1.075$ vs.\ $0.996$\,s, non-overlapping CIs but
a negligible absolute gap of $\approx\!80$\,ms) with p95 latency
within $\approx\!1$\,ms---the HSM decision layer is effectively free.
Against the scheduled policies, $\text{CETS}_{\mathrm{always}}$ is
$66.4$--$68.0\%$ and $\text{CETS}_{\mathrm{periodic},K=2}$ is
$44.0$--$51.8\%$: because no advisor call is economic, every trigger
the scheduled baselines issue is wasted work.  Trigger precision for
the scheduled policies is correspondingly low ($P_{\mathrm{strict}}
\!=\!0.100$ at periodic $K\!=\!2$, $0.095$ at always-on), since both
spend advisor cycles on windows that Theorem~\ref{thm:econ} proves
are sub-economic.

\emph{Receipt for Theorem~\ref{thm:econ}: mid-$Q$ selectivity.}
At $Q\!=\!700$ with $\thetastar\!=\!0.826$, HSM-gated fires exactly
once per run in every one of the 10 blocks, yielding
$P_{\mathrm{strict}}\!=\!1.000$ with a degenerate $[1.000,1.000]$
bootstrap interval, indicating that the gate both selects
the economically profitable transition
(\textsc{Burst\_Alt}\,$\to$\,\textsc{Burst\_Grp}, $S\!=\!0.781$,
$Q_i\!=\!557\!<\!700$) and correctly \emph{skips} the shallower one
($S\!=\!0.892$, $Q_1\!=\!1130\!>\!700$).  Versus always-on, HSM-gated
saves $66.5\%$ (Dexter) and $71.1\%$ (\texttt{index\_adv.}) of
wall-clock; versus periodic $K\!=\!2$, $50.7\%$ and $53.2\%$.
Precision is $100\times$--$200\times$ the scheduled baselines
($0.005$ always-on, $0.010$ periodic $K\!=\!2$).  Critically, p95
query latency under HSM-gated is \emph{equal to or better than} the
no-advisor baseline (Dexter: $42.59$ vs.\ $42.69$\,ms;
\texttt{index\_adv.}: $28.80$ vs.\ $42.69$\,ms with
non-overlapping CIs---Supabase's \texttt{index\_advisor} actually
reduces tail latency below baseline via its one economic rebuild).
This confirms HSM's skip decisions impose zero user-observable
penalty (\S\ref{sec:a13d} \emph{latency-on-skip} metric).

\emph{Practical $\theta$ bound: within-phase noise floor.}
Applying Theorem~\ref{thm:econ}'s $\thetastar(N,Q)\!=\!1\!-\!Q_{\min}/Q$
formula naively at $Q\!=\!1500$ yields $\thetastar\!=\!0.919$, which
over-fires every window because the within-phase signature in our
synthetic burst exhibits a \emph{noise floor} at
$S_{\mathrm{floor}}\!\approx\!0.875$ (random \texttt{aid}-sample
variability).  In general, the usable threshold is bounded below by
the stable-phase distribution:
$\thetastar_{\mathrm{eff}}(N,Q)\!=\!\min\!\bigl(\thetastar(N,Q),\,
S_{\mathrm{floor}}(W)\bigr)$, where $S_{\mathrm{floor}}(W)$ is the
$5$th-percentile within-phase similarity at window size $W$.  The
$Q\!=\!700$ operating point respects this bound
($0.826\!<\!0.875$); the $Q\!=\!1500$ operating point does not and
is therefore excluded from the main table.  For workloads whose
$Q_{\min}\!/(1\!-\!S_{\mathrm{floor}})\!>\!Q_i$ for every true
transition $i$, HSM-gated's economic recall is necessarily
$\le\!\lvert\{i:Q\!>\!Q_i\}\rvert\,/\,|\mathcal{T}|$---i.e.\ HSM
catches every transition that is economic \emph{and} separable from
within-phase drift.  Here that yields economic-recall $1/1$ at
$Q\!=\!700$ and strict-recall $1/2$ (the second transition being
correctly ignored per Theorem~\ref{thm:econ}).  We return to this in
\S\ref{sec:limitations}(iv).

\emph{B1 receipt: clamping prevents over-firing at large $Q$.}
At $Q\!=\!3000$, Theorem~\ref{thm:econ} predicts $\thetastar\!=\!0.959$, an aggressive bound.
The within-phase HSM signature exhibits a noise floor $S_{\mathrm{floor}}\!\approx\!0.875$ on this workload, so the clamp $\theta_{\mathrm{eff}}\!=\!\min(0.959,\,0.875)\!=\!0.875$ yields exactly $1$ trigger per block at precision $1.000$.
Without the clamp, the naive $\theta\!=\!0.959$ trigger fires $601$ times per block (a $601\times$ increase in advisor invocations) at precision $0.003$.
Wasted advisor cycles scale with advisor cost: $9.3\times$ more advisor time on Supabase \texttt{index\_advisor} ($4.66$\,s vs.\ $0.50$\,s) and $325\times$ more on Dexter ($58.5$\,s vs.\ $0.18$\,s).
Wall-clock impact depends on the advisor: Dexter degrades from $87$\,s to $146$\,s, while \texttt{index\_advisor} actually completes faster naive ($37$\,s vs.\ $61$\,s) because its lightweight recommend-only operations let early-built indexes amortise.
The headline value of clamping is therefore precision: it filters the $600$ false-positive triggers without trading on wall-clock for fast advisors.
The Supplementary (Cross-Workload $\theta^{*}$ Generalisability) separately confirms that the angular $S_R$ and $S_T$ dimensions are amplitude-invariant ($k\!\in\!\{0.5,1,2,5,10\}$), with HSM floor in $[0.695,0.875]$ across workloads (B4 receipt).

\FloatBarrier
\subsection{A-CE: Cross-Engine Validation on MongoDB}\label{sec:mongo-exp}

To validate HSM's DBMS-agnostic claim, we deploy the \emph{same} kernel
($W_0$, DWT+SAX+FastDTW) on MongoDB~7 with the document-store extractor.
A 24-window trace across four collection-specific phases (drift at W7, W13, W19)
uses the identical 10-block RCB protocol as PostgreSQL.
\textbf{A-CE1 ($\theta\!=\!0.75$):} DR\,=\,1.271; 61.7\,\% fewer advisor calls than always-on;
Youden $J\!=\!0.705$ (recall\,=\,1.000, specificity\,=\,0.705).

\textbf{Threshold Calibration.}
At $\theta\!=\!0.75$ (PostgreSQL parity), MongoDB shows $J\!=\!0.705$
because the polymorphic single-collection design yields $\bar{S}_V\!=\!1.0$
(no volume variation), lowering within-phase HSM to $\approx\!0.78$ vs.\
PostgreSQL's $\approx\!0.88$.
A post-hoc ROC sweep reveals
$\thetastar_{\mathrm{Mongo}}\!=\!0.55$ ($J^*\!=\!1.000$) vs.\
$\thetastar_{\mathrm{PG}}\!=\!0.775$ ($J^*\!=\!0.965$).

\textbf{A-CE2 ($\theta\!=\!0.65$, engine-calibrated):} Re-running the
\emph{identical} workload (same seeds) with $\theta\!=\!0.65$ yields
DR\,=\,1.486, $J\!=\!0.981$ (specificity 0.981, FP 4/210 vs.\ 62/210),
and 85.8\,\% fewer advisor calls---a $1.161\times$ speedup over always-on
(vs.\ $1.011\times$ at $\theta\!=\!0.75$).
This confirms that \emph{the HSM kernel is engine-agnostic; only $\theta$
requires per-deployment calibration}, consistent with the ROC framework
of Theorem~\ref{thm:gate}.

\emph{Cross-engine theta mismatch (B3 receipt).}
When the economically optimal $\theta$ from one engine is applied to a second engine without recalibration, precision degrades sharply.
B3-A: PostgreSQL at $\theta\!=\!0.55$ (MongoDB-optimal) yields zero triggers over 10 blocks;
B3-B: MongoDB at $\theta\!=\!0.775$ (PostgreSQL-optimal) yields $\approx\!10$ triggers per block with precision 0.27.
This validates Theorem~\ref{thm:gate}'s prediction that $\thetastar(N,Q)$ is workload- and engine-specific, requiring per-deployment calibration.

\textbf{A-CE3:} Per-window breakdown schema verified identical (16 columns) across engines.
The ROC curves and Youden-$J$ sweep are in Supplementary~\S XIV.

\FloatBarrier
\subsection{Summary of Experimental Results}

% v2 NOTE: Tables~\ref{tab:summary} and~\ref{tab:baselines} now reflect
% the HSM\_gated v2 throughput regeneration (SF\,$\in\{0.2,1.0,3.0\}$,
% 10 blocks, 4 conditions, 120 runs total; results/throughput.csv).
% Per-SF normalised retentions: HSM-gated 0.619/0.607/0.631; periodic
% 0.432/0.415/0.454; always-on 0.207/0.199/0.233.

Table~\ref{tab:summary} summarises all A1--A14 and A-CE outcomes with
their corresponding theorem validations. All six theorems and four
lemmas receive consistent empirical support on PostgreSQL; cross-engine
validation on MongoDB (A-CE) confirms the DBMS-agnostic property.

\begin{table}[!htb]\centering\scriptsize
\caption{Summary of experimental findings (A1--A14, A-CE, B1--B5). V.\ = validated against the indicated theorem or lemma. All six theorems and four lemmas receive consistent empirical support on PostgreSQL; B-series validates noise-floor properties and kernel/weight robustness.}\label{tab:summary}
\setlength{\tabcolsep}{3pt}
\begin{tabular}{@{}clp{4.8cm}c@{}}
\toprule
\textbf{RQ} & \textbf{Thm} & \textbf{Key result} & \textbf{V.}\\
\midrule
A1  & L1   & 4 axioms verified; 0 violations ($n\!=\!280$)       & $\checkmark$\\
A2  & L1   & $\SV\!\in\![0.35,1.00]$; unity for equal sizes      & $\checkmark$\\
A3  & L1   & $\ST$ within 0.988; cross 0.983                     & $\checkmark$\\
A4  & L1   & $\SA$ drops 63.2\% at phase boundaries              & $\checkmark$\\
A5  & L1   & DR\,=\,1.624 (95\%CI [1.613,1.636]); ICC\,=\,0.952; $p\!<\!0.001$; $|r|\!=\!0.999$ & $\checkmark$\\
A6  & T2,T3 & $\ST$ ($-$11.1\%), $\SA$ ($-$7.2\%) critical; speedup $8\times$  & $\checkmark$\\
A7  & L3   & OLS slope\,=\,0.935; $R^2\!=\!0.997$               & $\checkmark$\\
A8  & Gen. & SDSS 1.310; JOB 1.250; TPC-B 2.008; burst 1.093    & $\checkmark$\\
A9  & Util.& 0.25\% overhead; 919\,trans/s                       & $\checkmark$\\
A10 & Rob. & $\leq\!13$\% DR change at $\pm60$\% noise; $\leq\!4.1$\% at $\pm0.05$ weight & $\checkmark$\\
A11 & T3   & Break-even $\Npts\!\leq\!17{,}000$                  & $\checkmark$\\
A12 & T5   & $\beta\!=\!1.094$ ($R^2\!=\!0.9999$); $\Theta(N\log N)$ rebuild scaling ($N_{\mathrm{pts}}\!\in\![100,30000]$) & $\checkmark$\\
A13 & T4,T5 & HSM-gated $2.92\!\pm\!0.18\!\times$ vs.\ always-on, $1.43\!\pm\!0.04\!\times$ vs.\ periodic; $80.8\%$/$42.5\%$ fewer calls; $61.9\!\pm\!1.2\%$ no-advisor retention; scale-invariant CV\,$<\!7\%$ across SF\,$\in\{0.2,1.0,3.0\}$ & $\checkmark$\\
A13b & T4 & $z\!\ge\!3.40\sigma$ across scale; FAP$\!<\!1$\%; $d'\!>\!12$; Det.\,Prob\,=\,100\% & $\checkmark$\\
A13c & T5 & Bound $1/(1\!-\!p_{\mathrm{stable}})$ realised on TPC-H ($5.20\times$, $80.8\%/42.5\%$), burst ($13.0\times$, $80.8\%/44.4\%$), SDSS ($4.01\times$, $75.1\%/25.2\%$); OLTP correctly diagnosed as high-drift ($p_{\mathrm{stable}}\!\to\!0$, bound self-deactivates) & $\checkmark$\\
A14 & T4,T5 & Youden $J\!\ge\!0.9626$ across all SF; best at SF\,=\,0.2 ($J\!=\!0.9684$); 95\,\% certificate $[5.72\!\times\!,23.45\!\times]$ at $\pi_w\!=\!0.95$ (SF\,=\,3.0 worst) & $\checkmark$\\
A17 & T6   & Optimal $\Npts$ closed-form; empirical min at $\Npts\!=\!50$ (TPC-H SF\,0.2, SF\,3.0); $\Npts\!=\!100$ (JOB) & $\checkmark$\\
\midrule
A-CE1 & L1,T2 & MongoDB DR\,=\,1.271 ($\theta\!=\!0.75$); $J\!=\!0.705$ & $\checkmark$\\
A-CE2 & T4,T5 & $\theta\!=\!0.65$: DR\,=\,1.486; $J\!=\!0.981$; 85.8\,\% fewer calls  & $\checkmark$\\
A-CE3 & Repro.& Breakdown CSV schema: 16 columns, identical across engines & $\checkmark$\\
\midrule
B1 & T3 & Noise floor $S_{\mathrm{floor}}\!\approx\!0.875$; clamped $\theta\!=\!0.875$ yields 1 call/block, P=1.000; naive $\theta\!=\!0.959$ fires 601 times, P=0.003 & $\checkmark$\\
B2 & T2 & Kernel ablation: angular $J^*\!=\!0.667$ (AUC=0.833), cosine $J^*\!=\!0.667$ (AUC=0.833), euclidean $J^*\!=\!1.0$ (AUC=1.0) & $\checkmark$\\
B3 & T3 & Cross-engine mismatch: PG at $\theta\!=\!0.55$ (Mongo-opt) yields 0 triggers; Mongo at $\theta\!=\!0.775$ (PG-opt) yields $\approx$10 triggers, P=0.27 & $\checkmark$\\
B4 & T2 & Amplitude invariance: $S_R$, $S_T$ constant across $k\in\{0.5,1,2,5,10\}$; floor ranges 0.695--0.875 & $\checkmark$\\
B5 & T2 & Robustness: $\theta^*$ stable 0.882--0.920 (CV$<$1\%) across noise/weight perturbations; AUC degrades at $\delta\geq0.075$ & $\checkmark$\\
\bottomrule
\end{tabular}
\par\smallskip $\checkmark$\,=\,validated.
\end{table}

\FloatBarrier
\subsection{Discussion}

% v2 NOTE: A13/A12 magnitudes in this Discussion subsection have been
% re-checked against the HSM\_gated v2 sweep
% (results/throughput.csv; results/scale\_analysis.csv).  Speedup band,
% retention, scale-invariance CV, and $\beta$ now match the v2 numbers
% reported in Table~\ref{tab:throughput} and Figure~\ref{fig:throughput}.
% Other quantitative claims (Hoeffding certificate, weight-sensitivity)
% are sourced from the unchanged calibration / robustness analyses.

\textbf{Theory--Experiment Alignment.}
All theorems and lemmas receive consistent empirical support.
Scale-invariant DR (CV\,=\,0.1\%) confirms $\theta\!=\!0.75$ as robust.
Stratified Hoeffding yields $16.1\times$--$41.9\times$ central speedup at
$\pi_w\!\in\![0.95,0.99]$ (SF\,=\,3.0 inputs, worst case).
Inter-domain DR consistency (Supplementary, Dimensional-Utility Table)
supports Theorem~\ref{thm:minimal}; weight sensitivity (A10b) confirms
$\mathbf{w}_0$ lies in a broad flat optimum ($\leq5.5\%$ DR change).

\textbf{Threats to Validity.}
(i)~Phase leakage (N5) most influential ($\pm13$\%, A10).
(ii)~Three external validations (SDSS, JOB, pgbench) address TPC-H bias.
(iii)~OLTP may require weight recalibration; MongoDB (A-CE) tests data-model generality.
(iv)~J$\to$P/P$\to$C boundaries (HSM\,=\,0.77--0.83) reflect genuine query-mix similarity.

\textbf{Boundary Conditions.}
\textit{Gradual drift} may keep $\SHSM\!>\!\thetastar$ as indexes degrade;
operators should reduce $\Npts$ or adopt a sliding-window variant.
\textit{Concurrent phase changes} can cancel in the composite score;
per-dimension inspection provides diagnostic resolution.

\textbf{Baseline Comparison.}\label{sec:comparative}
On a single workload, the best single dimension can match HSM
(Table~\ref{tab:baselines}): $\SA$ alone reaches $J\!=\!0.993$ on
TPC-H, $\ST$ alone reaches $J\!=\!0.832$ on SDSS.  No single
dimension is consistently strong, however: $\SA$ drops to $0.227$
on SDSS and $\ST$ drops to $0.845$ on TPC-H.  The composite trades
peak per-workload $J$ for cross-workload consistency
($J\!\ge\!0.58$ on SDSS, $J\!\ge\!0.97$ on TPC-H), the formal
metric properties of Lemma~1, the dimension-wise diagnosis of
Algorithm~2, and the closed-form threshold of Theorems~3--5.
Operators with known workload characteristics may further improve
$J$ by tuning $\mathbf{w}$ toward the dominant dimension.

\textbf{Scope vs.\ ML-Based Advisors.}
BALANCE~\cite{Wang2024BALANCE} and Indexer++~\cite{Sharma2022IndexerPP}
solve the complementary \emph{what}-to-build problem; a throughput
comparison would conflate advisor quality with gate quality, and
public implementations are unavailable.

\begin{table}[!htb]\centering\scriptsize
\caption{Baseline comparison: Youden $J$ at optimal $\theta^{*}$
(all-pairs protocol). Bold: best single dimension per workload.
No single dimension achieves $J\!>\!0.85$ on all workloads;
only the composite HSM achieves consistent cross-workload detection
across all four benchmarks. Kernel-form rows (last block, marked $^{\dagger}$)
are evaluated on the burst-mid-$Q$ subset (kernel-ablation
figure in the Supplementary, A13 Decision-Layer Cost; B2 receipt)
and quantify pattern geometry rather than dimension utility.}\label{tab:baselines}
\setlength{\tabcolsep}{3pt}
\begin{tabular}{@{}lccc@{}}
\toprule
\textbf{Method} & \textbf{TPC-H} & \textbf{SDSS} & \textbf{MongoDB}\\
 & \textbf{SF\,0.2} & & \\
\midrule
$\SR$ only       & 0.979 & 0.416 & 1.000\\
$\SV$ only       & 0.000 & 0.001 & 0.000\\
$\ST$ only       & 0.849 & \textbf{0.832} & \textbf{1.000}\\
$\SA$ only       & \textbf{0.993} & 0.227 & 0.908\\
$\SPER$ only     & 0.022 & 0.086 & 0.985\\
\midrule
$\ST\!+\!\SA$ (2-dim) & 0.962 & 0.652 & 1.000\\
Uniform 5-dim    & 0.969 & 0.575 & 1.000\\
\midrule
Angular kernel$^{\dagger}$   & 0.667 & -- & --\\
Cosine kernel$^{\dagger}$    & 0.667 & -- & --\\
Euclidean dist.$^{\dagger}$  & 1.000 & -- & --\\
\midrule
\textbf{HSM ($\mathbf{w}_0$)} & \textbf{0.971} & \textbf{0.576} & \textbf{1.000}\\
\bottomrule
\end{tabular}
\end{table}

\textbf{A12/A14: Scale Sensitivity and Detector Quality
(Theorems~\ref{thm:gate} and~\ref{thm:speedup}).}
OLS regression of index-rebuild cost $T_A(N)$ across
$N_{\mathrm{pts}}\!\in\![100,30{,}000]$ yields
$\beta\!=\!1.094$ ($R^2\!=\!0.9999$), confirming $\Theta(N\log N)$.
Youden $J\!\ge\!0.9626$ at every scale; SF\,=\,0.2 achieves
$J^{*}\!=\!0.9684$ at $\thetastar\!=\!0.775$; SF\,=\,3.0 (worst case)
still achieves $J^{*}\!=\!0.9626$ at $\thetastar\!=\!0.775$.
Theorem~\ref{thm:speedup} gives a 95\,\% deployment certificate
$[5.72\times,23.45\times]$ at $\pi_w\!=\!0.95$ (SF\,=\,3.0 inputs,
worst case); central $16.1\times$ at $\pi_w\!=\!0.95$, rising to
$41.9\times$ at $\pi_w\!=\!0.99$.
Full per-SF tables, Hoeffding bands, ROC sweeps, and the scale
analysis figure are in Supplementary~\S XII.

\section{Conclusion}

HSM is a five-dimensional, training-free similarity measure for
proactive index management, with a formal core (six theorems, four
lemmas) covering metric soundness, $\Theta(\Npts)$ complexity, an
economic gating threshold, stratified-Hoeffding detector bounds, a
speedup certificate $\to\!1/(1{-}p_{\mathrm{stable}})$, and a
closed-form optimal window size.  Experiments confirm each theorem
(Table~\ref{tab:summary}): scale-invariant DR across SF\,0.2--3.0,
$<\!0.3\%$ overhead, and a central $16.1\times$ speedup at
$\pi_w\!=\!0.95$ (95\% interval $[5.72\times, 23.45\times]$ on the
worst-case SF\,=\,3.0 inputs).

\textit{Limitations.}\label{sec:limitations}
(i)~J$\to$P and P$\to$C boundaries (HSM\,=\,0.77--0.83) are not
flagged because the underlying query mixes are genuinely similar.
(ii)~Fixed weights $\mathbf{w}_0$ do not adapt to workload class:
on SDSS-like workloads dominated by $\ST$, a single-dimension
detector can outperform the composite
(Table~\ref{tab:baselines}); automated weight calibration would
mitigate this.
(iii)~End-to-end throughput is reported only for
SF\,$\in\{0.2,1.0,3.0\}$ (up to $\approx\!18$\,M rows on TPC-H).
The kernel-measurement cost itself is empirically
$\Theta(\Npts\log\Npts)$ on $\Npts\!\in\![100,\,30{,}000]$
($\beta\!=\!1.094$, $R^{2}\!=\!0.9999$), so the gating layer
itself extrapolates beyond this range; full end-to-end runs at
SF\,$\geq\!10$ are dominated by buffer-pool, autovacuum, and
unindexed-join effects outside the kernel and are deferred to
future work.
(iv)~Theorem~\ref{thm:speedup}'s operational speedup is
workload-specific by construction and asymptotes to $1\!\times$
on intrinsically high-drift workloads such as pgbench OLTP
($p_{\mathrm{stable}}\!\to\!0$); the bound remains exact, and the
gate's TPR\,=\,$1.00$ at default $\theta$ correctly reports the
diagnosis.
\textit{Future work:} live advisor integration; multi-tenant
cloud; automated $\theta$ calibration; extended NoSQL validation;
SF\,$\geq\!10$ end-to-end profiling.
Complete proofs are in Supplementary~\S I--X.

\section*{Code and Data Availability}
A complete replication package---source code, configuration, scripts, and
instructions to reproduce all reported results---is available from the
author upon reasonable request during review, and will be publicly released
at a citable archive upon acceptance.

\section*{Acknowledgments}
This work used generative AI tools per IEEE guidelines:
Claude (Anthropic) for code review, formula verification, and language checking;
SciSpace for literature search and writing polish.

\bibliographystyle{plain}
\bibliography{references}
\balance
\end{document}
